\documentclass[11pt,a4paper]{article}
\usepackage[T1]{fontenc}
\usepackage[utf8]{inputenc}
\usepackage{lmodern}
\usepackage{microtype}
\usepackage[margin=25mm,headheight=15pt]{geometry}
\usepackage{amsmath,amssymb,amsthm,mathtools}
\usepackage{booktabs,longtable,array,enumitem,xcolor}
\usepackage{fancyhdr,titlesec,multicol}
\usepackage[most]{tcolorbox}
\usepackage{xurl}
\usepackage{hyperref}
\definecolor{ink}{HTML}{142A43}
\definecolor{linkblue}{HTML}{234F75}
\definecolor{pale}{HTML}{F1F5F8}
\hypersetup{colorlinks=true,linkcolor=linkblue,citecolor=linkblue,urlcolor=linkblue,
  pdftitle={Hilbert-Vector Resolvents and the Riemann Hypothesis},
  pdfauthor={Riemann-Hypothesis-Final project research note},
  pdfsubject={Detailed on-paper analytic seed criterion; no unconditional RH or Lean certification},
  pdfkeywords={Riemann hypothesis, xi function, Hilbert space, resolvent, Hyperbolic Zeta, positive determinant, finite witness}}
\urlstyle{same}
\pagestyle{fancy}
\fancyhf{}
\fancyhead[L]{\small\textsc{Hilbert-vector RH anchor}}
\fancyhead[R]{\small Analytic construction}
\fancyfoot[L]{\footnotesize 27 October 2026 \enspace | \enspace v1.1}
\fancyfoot[R]{\small\thepage}
\renewcommand{\headrulewidth}{0.3pt}
\titleformat{\section}{\large\bfseries\color{ink}}{\thesection}{0.65em}{}
\titleformat{\subsection}{\normalsize\bfseries\color{ink}}{\thesubsection}{0.65em}{}
\setlength{\parskip}{4pt}
\setlength{\parindent}{0pt}
\setlist{topsep=4pt,itemsep=3pt,parsep=0pt}
\setcounter{tocdepth}{1}
\numberwithin{equation}{section}
\newtheorem{theorem}{Theorem}[section]
\newtheorem{proposition}[theorem]{Proposition}
\newtheorem{lemma}[theorem]{Lemma}
\newtheorem{corollary}[theorem]{Corollary}
\theoremstyle{definition}
\newtheorem{definition}[theorem]{Definition}
\newtheorem{example}[theorem]{Example}
\theoremstyle{remark}
\newtheorem{remark}[theorem]{Remark}
\newcommand{\C}{\mathbb C}
\newcommand{\R}{\mathbb R}
\newcommand{\N}{\mathbb N}
\newcommand{\HH}{\mathcal H}
\newcommand{\BB}{\mathcal B}
\newcommand{\id}{I}
\newcommand{\ip}[2]{\langle #1,#2\rangle}
\newcommand{\norm}[1]{\lVert #1\rVert}
\newcommand{\RH}{\mathrm{RH}}
\newcommand{\HSeed}{\mathsf{HSeed}_{\Theta}}
\newcommand{\DSeed}{\mathsf{DSeed}_{\Theta}}
\newcommand{\MSeed}{\mathsf{MSeed}_{\Theta}}
\newcommand{\dd}{\,\mathrm d}
\DeclareMathOperator{\Tr}{Tr}
\DeclareMathOperator{\supp}{supp}
\DeclareMathOperator{\Repart}{Re}
\DeclareMathOperator{\Impart}{Im}
\newtcolorbox{scopebox}{colback=pale,colframe=ink!35,boxrule=0.4pt,arc=0pt,
  left=9pt,right=9pt,top=6pt,bottom=6pt,breakable}
\newcommand{\reposha}{\texttt{362d9f8aeaaf9479cc27ceb6b63bdb7456267990}}
\begin{document}
\begin{center}
{\small\sffamily\color{linkblue} SOLOBOLO \enspace / \enspace RIEMANN-HYPOTHESIS-FINAL}\par
\vspace{9pt}
{\LARGE\bfseries\color{ink} Hilbert-Vector Resolvents\par
and the Riemann Hypothesis}\par
\vspace{7pt}
{\large A detailed analytic seed criterion}\par
\vspace{9pt}
{\small Project research note \enspace | \enspace 27 October 2026}\par
{\small Version 1.1 \enspace | \enspace B-Tier derivation \enspace | \enspace construction}
\end{center}
\vspace{5pt}
\begin{abstract}
We give a precise, self-contained analytic formulation of a Hilbert-vector
anchor for the Riemann hypothesis. Let $\Theta$ be the entire function
determined by $\Theta(w^2)=\xi(\tfrac12+w)$, and put
$L(a)=\Theta'(a)/\Theta(a)$ for $a\geq0$. We prove that RH is equivalent
to the existence of one Hilbert space, one bounded nonnegative self-adjoint
operator $K$, and one fixed vector $v$ such that
$\frac{\mathrm d}{\mathrm da}(aL(a))=\norm{(I+aK)^{-1}v}^2$
on a nonempty positive interval. The derivation includes the resolvent
calculus, its divided-difference Gram kernel, and the analytic continuation
that excludes zeros off the negative real axis. A possible $C/a$ ambiguity
is removed using $\Theta(0)\ne0$, rather than suppressed by an unstated
normalization. The reverse implication constructs a positive diagonal
operator from the zero product under RH, with multiplicities counted
correctly. We also develop the spectral-measure and all-order moment
formulations, explain basepoint shifts, and prove that no exactly calibrated
finite-dimensional model can suffice. The result is an analytic equivalence;
it does not construct the required seed unconditionally or produce an
accepted finite witness. Hyperbolic zeta identities additionally give
explicit comparison models and a conditional determinant-limit route
to the actual seed.
\end{abstract}
\begin{scopebox}
\small\textbf{Evidence and scope.} Every proof in this note is an on-paper
derivation. No Lean compiler or axiom audit was run. The seed predicates
below are defined explicitly here; they are not asserted to be recovered
field-for-field definitions of a legacy \texttt{HasSeed} declaration.
The repository's finite-checker target remains open. Classical background
is cited, and no novelty or unconditional proof of RH is claimed.
\end{scopebox}
\vspace{5pt}
\section*{Contents}
\begin{multicols}{2}
\small\setlength{\parskip}{0pt}
\makeatletter\@starttoc{toc}\makeatother
\end{multicols}
\clearpage

\section[Assertion and scope]{The exact assertion and its role}\label{sec:assertion}

The starting point is the family of identities
\begin{equation}\label{eq:anchor-intro}
R_a=(I+aK)^{-1},\qquad
\frac{\mathrm d}{\mathrm da}\bigl(aL(a)\bigr)=\norm{R_av}^{2},\qquad
V_a=2\sqrt a\,R_av.
\end{equation}
Their significance depends on their quantifiers. The operator and vector
must be the \emph{same} for every parameter in the interval, and the scalar
$L$ must come from the actual completed Riemann zeta function. An arbitrary
nonnegative scalar function can be represented as a squared norm by
allowing its vector to vary freely. Such a pointwise representation contains
none of the resolvent compatibility used below.

We fix the standard completion
\begin{equation}\label{eq:xi}
\xi(s)=\tfrac12s(s-1)\pi^{-s/2}\Gamma(s/2)\zeta(s),
\end{equation}
with the removable values understood. The symmetry $\xi(s)=\xi(1-s)$
and normalization are classical \cite{dlmf-reflection}. Section~\ref{sec:theta}
defines $\Theta$ without choosing a square-root branch and proves that
$L$ is a well-defined real analytic function on $[0,\infty)$.

\begin{definition}[Calibrated Hilbert seed]\label{def:hseed}
The assertion $\HSeed$ means that there exist a complex Hilbert space
$\HH$, an operator $K\in\BB(\HH)$, a vector $v\in\HH$, and real numbers
$0<\alpha<\beta<\infty$ such that
\begin{equation}\label{eq:seed-hypotheses}
K=K^*,\qquad \ip{x}{Kx}\geq0\quad(x\in\HH),
\end{equation}
and, for every $a\in(\alpha,\beta)$,
\begin{equation}\label{eq:hseed}
L(a)=\ip{v}{(I+aK)^{-1}v}.
\end{equation}
The inner product is conjugate-linear in its first slot and linear in its
second slot. In particular, the complex extension
$z\mapsto\ip{v}{(I+zK)^{-1}v}$ is holomorphic where the inverse exists.
\end{definition}

\begin{definition}[Differential Hilbert seed]\label{def:dseed}
The assertion $\DSeed$ has the same existential data and operator
hypotheses as Definition~\ref{def:hseed}, but replaces
\eqref{eq:hseed} by
\begin{equation}\label{eq:dseed}
\frac{\mathrm d}{\mathrm da}\bigl(aL(a)\bigr)
=\norm{(I+aK)^{-1}v}^{2}
\quad\text{for every }a\in(\alpha,\beta).
\end{equation}
This is an exact functional identity on an interval, not a test at finitely
many sample points.
\end{definition}

\begin{theorem}[Main analytic criterion]\label{thm:main}
For the actual $\Theta$ associated with \eqref{eq:xi},
\begin{equation}\label{eq:main}
\RH\quad\Longleftrightarrow\quad\HSeed
\quad\Longleftrightarrow\quad\DSeed.
\end{equation}
If either seed identity holds on one nonempty positive interval, its
calibration extends to every $a\geq0$ for the same $K$ and $v$.
Under RH, a seed can be chosen on $\ell^2(\N;\C)$ with $K$ positive,
compact and trace class. In that realization,
\begin{equation}\label{eq:normalization-intro}
\Tr K=\norm v^2=\frac{\xi''(1/2)}{2\xi(1/2)}.
\end{equation}
\end{theorem}

Sections~\ref{sec:resolvent}--\ref{sec:forward} prove the forward and
differential implications without assuming RH or a positive zero product.
Section~\ref{sec:converse} proves the converse under RH. The theorem is
a criterion: the existential seed assertion carries the mathematical
difficulty of RH. No independently constructed seed is supplied here.

There are two different questions behind this note. The analytic question
asks exactly what a fixed positive resolvent identity implies. The finite
question asks whether a terminating checker on finite raw data can certify
such an identity, and whether RH guarantees an accepted input to that same
checker. The second question does not follow merely by naming the first
one a seed. Section~\ref{sec:finite} makes this distinction precise.

\section[Entire function and zero geometry]{The actual entire function and its zero geometry}\label{sec:theta}

\subsection{Square descent without a branch choice}

Put $E(w)=\xi(\tfrac12+w)$. This is entire, and the functional equation
makes it even. Its Taylor series has the form
\begin{equation}\label{eq:theta-series}
E(w)=\sum_{n=0}^{\infty}c_nw^{2n},\qquad
c_n=\frac{\xi^{(2n)}(1/2)}{(2n)!}.
\end{equation}
Define
\begin{equation}\label{eq:theta-definition}
\Theta(z)=\sum_{n=0}^{\infty}c_nz^n.
\end{equation}
The series has infinite radius of convergence. Indeed, for any $r>0$,
the absolute convergence of \eqref{eq:theta-series} at a point with
$|w|=\sqrt r$ gives $\sum_n|c_n|r^n<\infty$.
Consequently
\begin{equation}\label{eq:square-descent}
\Theta(w^2)=\xi(\tfrac12+w)\qquad(w\in\C).
\end{equation}
The shorthand $\Theta(z)=\xi(\tfrac12+\sqrt z)$ therefore has a
single-valued meaning: the two square roots give the same value. The branch
cut introduced later belongs to the resolvent domain, not to $\Theta$.

Complex conjugation preserves $\xi$, as follows first on its defining
real domain and then by analytic continuation. Thus the coefficients
$c_n$ are real, and
$\Theta(\overline z)=\overline{\Theta(z)}$.

\subsection{Positivity on the real positive axis}

For completeness, the nonvanishing needed to define $L$ can be checked
without RH. For $0<s<1$, the alternating-series formula
\cite[Eq.~25.2.3]{dlmf-series} gives
\begin{equation}\label{eq:eta}
\zeta(s)=\frac{\eta(s)}{1-2^{1-s}},\qquad
\eta(s)=\sum_{n=1}^{\infty}(-1)^{n-1}n^{-s}>0.
\end{equation}
The strict inequality follows by grouping consecutive terms:
$(2j-1)^{-s}-(2j)^{-s}>0$. The even partial sums converge to the same
alternating-series limit and already their first pair is positive.
The denominator in \eqref{eq:eta} is negative, so $\zeta(s)<0$.
Since $s(s-1)<0$ and the remaining factors in \eqref{eq:xi} are positive,
$\xi(s)>0$ on $(0,1)$. For $s>1$, every factor in \eqref{eq:xi}
is positive, and the removable value at $s=1$ is $1/2$.
Applying this with $s=1/2+\sqrt a$ proves
\begin{equation}\label{eq:theta-positive}
\Theta(a)>0\quad(a\geq0),\qquad\Theta(0)=\xi(1/2)>0.
\end{equation}

It follows that
\begin{equation}\label{eq:L}
L(a)=\frac{\Theta'(a)}{\Theta(a)}\in\R\qquad(a\geq0)
\end{equation}
is real analytic, including at zero. On this axis it equals the recovered
quantity $\Repart(\Theta'(a)/\Theta(a))$; taking the real part loses no
information there. It would not be legitimate to make that replacement
for arbitrary complex $z$.

Expanding at zero gives the exact normalization
\begin{equation}\label{eq:kappa}
\kappa:=L(0)=\frac{\Theta'(0)}{\Theta(0)}
=\frac{\xi''(1/2)}{2\xi(1/2)}.
\end{equation}
No numerical approximation to $\kappa$ is needed in any proof below.

\subsection{The critical line becomes the negative real axis}

The zeros of $\xi$ are exactly the nontrivial zeros of $\zeta$; the
standard zero symmetry and the meaning of RH are recalled in
\cite{dlmf-zeros}. If $\rho=1/2+w$ is such a zero, then
$z=w^2$ is a zero of $\Theta$, and conversely every zero of $\Theta$
arises this way by \eqref{eq:square-descent}.

For $w=x+iy$,
\begin{equation}\label{eq:square-real}
w^2=(x^2-y^2)+2ixy.
\end{equation}
If $w^2$ is real and nonpositive, then $xy=0$ and $x^2-y^2\leq0$.
If $y=0$, these imply $x=0$; if $y\ne0$, the equality $xy=0$ again
implies $x=0$. Conversely $x=0$ gives $w^2=-y^2\leq0$.
Since $\Theta(0)\ne0$, its zeros never correspond to $w=0$.
We have therefore proved
\begin{proposition}[Exact zero geometry]\label{prop:geometry}
\begin{equation}\label{eq:rh-geometry}
\RH\quad\Longleftrightarrow\quad
\bigl(\Theta(z)=0\ \Longrightarrow\ z\in(-\infty,0)\bigr).
\end{equation}
\end{proposition}
This elementary change of variables is the final geometric step of the
forward proof. It refers to the actual $\xi$ and $\Theta$ throughout.

\section[Fixed-vector resolvent calculus]{Resolvent calculus with one fixed vector}\label{sec:resolvent}

Fix a complex Hilbert space $\HH$, a bounded nonnegative self-adjoint
operator $K$, and a vector $v$. For $a\geq0$, set
\begin{equation}\label{eq:Rfg}
R_a=(I+aK)^{-1},\qquad f(a)=\ip v{R_av},\qquad g(a)=af(a).
\end{equation}
At this stage $f$ is a model scalar function; equality with the actual
$L$ has not been assumed.

\subsection{Invertibility, self-adjointness and commutation}

For every $x\in\HH$, positivity and Cauchy--Schwarz imply
\[
\norm x^2\leq\ip x{(I+aK)x}\leq\norm x\,\norm{(I+aK)x}.
\]
Thus $I+aK$ is bounded below by $1$. Its range is closed. Its orthogonal
complement is the kernel of its adjoint, which equals its own kernel
because the operator is self-adjoint. The lower bound makes that kernel
zero, so the range is also dense and hence is all of $\HH$.
It follows that
\begin{equation}\label{eq:Rbounds}
\norm{R_a}\leq1,\qquad R_a^*=R_a,\qquad R_a\geq0.
\end{equation}
For the last inequality write $y=(I+aK)x$ and observe that
$\ip y{R_ay}=\ip{(I+aK)x}x\geq0$.
The operator $K$ commutes with $I+aK$ and therefore with its inverse.
The operators $R_a$ and $R_b$ also commute.

\begin{lemma}[Two resolvent identities]\label{lem:resolvents}
For $a,b\geq0$,
\begin{align}
R_a-R_b&=(b-a)R_aKR_b,\label{eq:first-resolvent}\\
aR_a-bR_b&=(a-b)R_aR_b.\label{eq:second-resolvent}
\end{align}
\end{lemma}
\begin{proof}
For the first identity, multiply out
$R_a[(I+bK)-(I+aK)]R_b$.
For the second, the same calculation with
$a(I+bK)-b(I+aK)=(a-b)I$ gives the result.
Both formulas also hold when $a=b$, with both sides then zero.
\end{proof}

\subsection{Differentiation in operator norm}

From \eqref{eq:first-resolvent},
\[
\frac{R_{a+h}-R_a}{h}=-R_{a+h}KR_a.
\]
Continuity of inversion, or the same identity and local boundedness, shows
that the right side tends in operator norm to $-R_aKR_a$.
Since $K$ commutes with $R_a$,
\begin{equation}\label{eq:Rprime}
R'_a=-KR_a^2.
\end{equation}
Repeated differentiation yields
\begin{equation}\label{eq:Rderivatives}
R_a^{(n)}=(-1)^n n!K^nR_a^{n+1}\qquad(n\geq0).
\end{equation}
One can verify the induction directly: differentiating $R_a^{n+1}$
produces $n+1$ identical commuting terms
$-KR_a^{n+2}$.
Applying the continuous linear functional $T\mapsto\ip v{Tv}$ gives
\begin{equation}\label{eq:fderivatives}
(-1)^n f^{(n)}(a)=n!\ip v{K^nR_a^{n+1}v}\geq0.
\end{equation}
The last operator is positive, for example by the bounded self-adjoint
functional calculus \cite{williams}. Thus $f$ is completely monotone.
This consequence alone will not characterize a calibrated seed.

\subsection{The norm-square identity}

The identity $(I+aK)R_a=I$ implies
$R_a=R_a^2+aKR_a^2$. Consequently
\begin{align}
g'(a)&=f(a)+af'(a)\notag\\
&=\ip v{(R_a-aKR_a^2)v}\notag\\
&=\ip v{R_a^2v}=\norm{R_av}^2.\label{eq:norm-square}
\end{align}
This proves the central algebraic formula, including its exact coefficient.
It involves $af(a)$, not $f(a)$ alone and not an unspecified change of
basepoint. Differentiating further gives
\begin{equation}\label{eq:gderivatives}
(-1)^n g^{(n+1)}(a)
=(n+1)!\ip v{K^nR_a^{n+2}v}\geq0\qquad(n\geq0).
\end{equation}
In particular $g$ is increasing and concave on the positive axis.
If $v\ne0$, then $R_av\ne0$ and $g'(a)>0$.

\section[The Gram kernel]{The Gram kernel and its exact normalization}\label{sec:gram}

Define $V_a=2\sqrt a\,R_av$ for $a>0$.
Taking the scalar product of \eqref{eq:second-resolvent} with $v$
and using self-adjointness gives, for $a\ne b$,
\begin{equation}\label{eq:unscaled-gram}
\frac{g(a)-g(b)}{a-b}
=\ip v{R_aR_bv}=\ip{R_av}{R_bv}.
\end{equation}
Hence
\begin{equation}\label{eq:gram}
\ip{V_a}{V_b}
=4\sqrt{ab}\,\frac{af(a)-bf(b)}{a-b}\qquad(a,b>0,\ a\ne b).
\end{equation}
If the calibration $f=L$ holds, this is precisely the recovered
Hilbert-vector Gram identity.

\begin{proposition}[Diagonal and positivity]\label{prop:gram}
The divided-difference kernel
\begin{equation}\label{eq:kernel}
\mathcal K_g(a,b)=
\begin{cases}
\dfrac{g(a)-g(b)}{a-b},&a\ne b,\\[5pt]
g'(a),&a=b
\end{cases}
\end{equation}
is positive semidefinite. More precisely, for any finite collection
$a_1,\ldots,a_m>0$ and complex numbers $c_1,\ldots,c_m$,
\begin{equation}\label{eq:gram-psd}
\sum_{i,j=1}^m\overline{c_i}c_j\mathcal K_g(a_i,a_j)
=\norm{\sum_{j=1}^m c_jR_{a_j}v}^2\geq0.
\end{equation}
The diagonal extension of \eqref{eq:gram} is
\begin{equation}\label{eq:gram-diagonal}
\norm{V_a}^2=4a\,g'(a)=4a\,\norm{R_av}^2.
\end{equation}
\end{proposition}
\begin{proof}
Equation~\eqref{eq:unscaled-gram} gives all off-diagonal entries, while
\eqref{eq:norm-square} gives the diagonal. Expanding the displayed squared
norm proves \eqref{eq:gram-psd}; multiplication by the factors
$2\sqrt{a_i}$ gives the scaled Gram matrix.
\end{proof}

For two distinct parameters, the determinant condition is the concrete inequality
\begin{equation}\label{eq:two-point}
\left|\frac{g(a)-g(b)}{a-b}\right|^2\leq g'(a)g'(b).
\end{equation}
It is a necessary condition tying different points together. Positivity of
$g'$ at every individual point does not imply it.

\begin{example}[Even strong scalar positivity is insufficient]\label{ex:scalar}
Let $\widetilde g(a)=1-e^{-a}$ and
$\widetilde f(a)=(1-e^{-a})/a$ for $a>0$.
Then $\widetilde g'>0$, $\widetilde g''<0$, and
\[
\widetilde f(a)=\int_0^1 e^{-at}\dd t,\qquad
(-1)^n\widetilde f^{(n)}(a)=\int_0^1t^ne^{-at}\dd t>0.
\]
Nevertheless, writing $d=a-b\ne0$, the off-diagonal divided difference is
\[
\mathcal K_{\widetilde g}(a,b)
=e^{-(a+b)/2}\frac{\sinh(d/2)}{d/2}.
\]
The determinant of the corresponding two-by-two matrix equals
\begin{equation}\label{eq:counterexample-det}
e^{-a-b}\left[1-
\left(\frac{\sinh(d/2)}{d/2}\right)^2\right]<0.
\end{equation}
The strict inequality follows from $|\sinh x|>|x|$ for nonzero real $x$,
which follows by integrating $\cosh t>1$ away from zero.
Thus these scalar monotonicity conditions cannot be substituted for the
fixed-resolvent structure.
\end{example}

\section[Holomorphic model]{Holomorphic extension of the model}\label{sec:analytic}

Choose a real number $M>\norm K$ with $M>0$, and set
\begin{equation}\label{eq:domain}
D_M=\C\setminus(-\infty,-1/M].
\end{equation}
This open connected domain contains zero and the entire positive real
axis. The bounded self-adjoint spectral theorem puts $\sigma(K)$ in
$[0,M]$ \cite{williams}. For $z\in D_M$ and $t\in[0,M]$, the number
$1+zt$ is nonzero. Indeed, a zero would require $t>0$ and
$z=-1/t\leq-1/M$. The functional calculus therefore defines
\begin{equation}\label{eq:complex-model}
R_z=(I+zK)^{-1},\qquad F(z)=\ip v{R_zv}\qquad(z\in D_M).
\end{equation}

\begin{lemma}[Norm-holomorphic resolvent]\label{lem:holomorphic}
The maps $z\mapsto R_z$ and $z\mapsto F(z)$ are holomorphic on $D_M$.
\end{lemma}
\begin{proof}
Fix $z_0\in D_M$. For sufficiently small $h$,
$\norm{hKR_{z_0}}<1$ and
\begin{align*}
I+(z_0+h)K&=(I+z_0K)(I+hKR_{z_0}),\\
R_{z_0+h}&=R_{z_0}\sum_{n=0}^{\infty}(-hKR_{z_0})^n.
\end{align*}
The series converges in operator norm on a neighborhood of $h=0$, so
it is a norm-convergent local power series. The scalar functional
$T\mapsto\ip v{Tv}$ is continuous and complex linear under the convention
in Definition~\ref{def:hseed}, and applying it proves the scalar assertion.
\end{proof}

It is useful to retain zero in $D_M$. A cut at $(-\infty,0]$ would still
suffice for the normalized seed-to-RH argument, but it would hide the
evaluation at zero that removes the integration constant in the
differential seed. Boundedness of $K$ supplies a genuine neighborhood of
zero on which the resolvent is analytic.

The spectral theorem also gives a scalar measure. If $E_K$ is the
projection-valued spectral measure, define
\begin{equation}\label{eq:measure}
\mu(B)=\ip v{E_K(B)v}\qquad(B\text{ Borel in }[0,M]).
\end{equation}
It is a finite positive measure, with total mass $\norm v^2$, and
\begin{equation}\label{eq:measure-resolvent}
F(z)=\int_{[0,M]}\frac{1}{1+zt}\dd\mu(t).
\end{equation}
For a compact $C\subset D_M$, the continuous function $|1+zt|$ has a
strictly positive minimum on $C\times[0,M]$. This bounds every fixed
derivative of the integrand by an integrable constant. Thus
\begin{equation}\label{eq:measure-derivatives}
F^{(n)}(z)=(-1)^n n!\int_{[0,M]}
\frac{t^n}{(1+zt)^{n+1}}\dd\mu(t).
\end{equation}
This justifies the integral differentiations directly and independently
checks the operator formulas of Section~\ref{sec:resolvent}.

\section[Transfer to RH]{The analytic transfer to RH}\label{sec:forward}

\subsection{Exact calibration on an interval is enough}

\begin{theorem}[Normalized seed implies RH]\label{thm:forward}
If $\HSeed$ holds, then $\RH$ holds. For its fixed seed data,
$F(a)=L(a)$ for all $a\geq0$.
\end{theorem}
\begin{proof}
On the interval in Definition~\ref{def:hseed}, the calibration and
$\Theta(a)>0$ imply
\[
\Theta'(a)-F(a)\Theta(a)=0.
\]
The function
\begin{equation}\label{eq:division-free}
G(z)=\Theta'(z)-F(z)\Theta(z)
\end{equation}
is holomorphic on $D_M$. Its zeros have an accumulation point there,
so the identity theorem gives
\begin{equation}\label{eq:ode}
\Theta'(z)=F(z)\Theta(z)\qquad(z\in D_M).
\end{equation}
This passage does not assume that $\Theta$ is zero-free on $D_M$.
The expression \eqref{eq:division-free} has no division by $\Theta$.

Suppose, for a contradiction, that $\Theta$ has a zero $z_0\in D_M$.
Since $\Theta(0)>0$, it is not identically zero. Its zero at $z_0$ has
some finite order $m\geq1$, so near $z_0$,
\[
\Theta(z)=(z-z_0)^m h(z),\qquad h(z_0)\ne0,
\]
where $h$ is holomorphic. Substituting in \eqref{eq:ode} and cancelling
$(z-z_0)^{m-1}$ off $z_0$ gives
\[
mh(z)+(z-z_0)h'(z)=(z-z_0)F(z)h(z).
\]
Both sides are continuous at $z_0$; evaluation there yields
$mh(z_0)=0$, a contradiction.
Therefore every zero of $\Theta$ lies outside $D_M$, hence in
$(-\infty,-1/M]\subset(-\infty,0)$.
Proposition~\ref{prop:geometry} gives RH. Finally, restricting
\eqref{eq:ode} to $a\geq0$ and dividing by the already known positive
$\Theta(a)$ proves $F(a)=L(a)$ everywhere on that ray.
\end{proof}

The proof uses neither a preselected logarithm of $\Theta$ on the cut
plane nor an RH-dependent product. Choosing a global logarithm before
proving zero-freeness would conceal the central difficulty. The
division-free ODE avoids that issue.

\subsection{The integration constant is real and must be treated}

Suppose only the differential identity \eqref{eq:dseed} holds. From
\eqref{eq:norm-square},
\[
\frac{\mathrm d}{\mathrm da}\bigl(aL(a)-af(a)\bigr)=0
\quad(a\in(\alpha,\beta)).
\]
The ordinary mean value theorem on subintervals shows that
\begin{equation}\label{eq:gauge}
a\bigl(L(a)-F(a)\bigr)=C,
\qquad L(a)=F(a)+\frac Ca
\end{equation}
for a constant $C\in\R$. Neither the derivative identity nor the
off-diagonal Gram identity alone detects this constant: adding $C$ to
$aL(a)$ changes neither its derivative nor its divided differences.

For an arbitrary function on an interval bounded away from zero, one
would need a normalization value to set $C=0$. The actual entire
$\Theta$ supplies a stronger analytic normalization.

\begin{theorem}[Differential seed and normalized seed agree]\label{thm:normalization}
For the actual $\Theta$, $\DSeed\Longleftrightarrow\HSeed$.
The same operator and vector can be used in both directions.
\end{theorem}
\begin{proof}
If $\HSeed$ holds, differentiating $aL(a)=af(a)$ and using
\eqref{eq:norm-square} gives $\DSeed$.
Conversely assume $\DSeed$ and let $C$ be the constant in
\eqref{eq:gauge}. Consider
\begin{equation}\label{eq:Q}
Q(z)=z\Theta'(z)-\bigl(zF(z)+C\bigr)\Theta(z).
\end{equation}
This is holomorphic on all of $D_M$, including zero. By
\eqref{eq:gauge}, it vanishes on $(\alpha,\beta)$, so the identity
theorem gives $Q=0$ throughout $D_M$. At zero,
\[
0=Q(0)=-C\Theta(0).
\]
Equation~\eqref{eq:theta-positive} therefore forces $C=0$.
It follows that $z(\Theta'-F\Theta)=0$ on $D_M$. Away from zero
we may divide by $z$, and continuity supplies the equality at zero.
Thus \eqref{eq:ode} holds, and positive-ray division by $\Theta(a)$
gives the normalized calibration. In particular $\HSeed$ holds.
\end{proof}

\begin{corollary}[The recovered norm-square identity is an RH criterion]
For a fixed bounded nonnegative self-adjoint $K$ and fixed $v$,
the identity \eqref{eq:dseed} on any nonempty positive interval
implies RH. No extra origin-normalization premise is needed beyond
the established facts about the actual $\Theta$.
\end{corollary}

The qualifier ``actual'' is essential. If one replaced $\Theta$ by
$z^m\Theta$ with $m>0$, its logarithmic derivative would acquire
$m/z$, and the derivative of $aL$ would not see that origin factor.
The known value $\Theta(0)\ne0$ is precisely what rules out this
ambiguity here.

\section[Converse under RH]{The converse: constructing a seed under RH}\label{sec:converse}

\subsection{Why the product has no exponential ambiguity}

We use two classical facts about entire functions. The completed
$\xi$ has order at most $1$; a direct growth argument from the theta
integral is given in \cite{garrett}. Hadamard factorization and its
zero-summability consequence are stated in
\cite[Theorems~3.7 and 3.11]{mcmullen}.
To see exactly how these apply, write
$M_E(r)=\max_{|w|=r}|E(w)|$ and define $M_\Theta$ similarly.
The map $w\mapsto w^2$ maps the circle of radius $\sqrt r$ onto the
circle of radius $r$, so
\begin{equation}\label{eq:order-halving}
M_\Theta(r)=M_E(\sqrt r).
\end{equation}
A bound $|E(w)|\leq C_\varepsilon
\exp(C_\varepsilon |w|^{1+\varepsilon})$ therefore puts the order
of $\Theta$ at most $1/2$. Only the strict inequality ``order less
than $1$'' is needed below.

For an entire function of order less than $1$, nonzero at zero,
Hadamard factorization has genus zero and a constant exponential factor.
Its nonzero zeros $z_n$, repeated by multiplicity, satisfy
$\sum_n|z_n|^{-1}<\infty$. Hence
\begin{equation}\label{eq:general-product}
\Theta(z)=\Theta(0)\prod_n\left(1-\frac z{z_n}\right).
\end{equation}
There is no factor $e^{bz}$ with $b\ne0$: the polynomial in the
Hadamard exponential has degree at most the order, hence degree zero.
There is also no origin monomial because $\Theta(0)\ne0$.

\subsection{Zero pairs and multiplicity}

Assume RH. Enumerate the positive ordinates of the nontrivial zeros
as $\gamma_1,\gamma_2,\ldots$, repeating each according to its
multiplicity. There are infinitely many zeros, with no finite
accumulation point \cite{dlmf-zeros}. We can arrange
$0<\gamma_1\leq\gamma_2\leq\cdots$ with $\gamma_n\to\infty$.
The two points $w=\pm i\gamma$ are the preimages of
$z=-\gamma^2$. Since the derivative of $w\mapsto w^2$ is nonzero
at either preimage, a zero of multiplicity $m$ at either point
corresponds to a zero of multiplicity $m$ of $\Theta$.
It does \emph{not} become a zero of multiplicity $2m$.
Equation~\eqref{eq:general-product} becomes
\begin{equation}\label{eq:rh-product}
\Theta(z)=\Theta(0)\prod_{n=1}^{\infty}
\left(1+\frac z{\gamma_n^2}\right),\qquad
\sum_{n=1}^{\infty}\frac1{\gamma_n^2}<\infty.
\end{equation}
Thus one positive/negative pair contributes one factor for each
occurrence of its multiplicity.

\subsection{Logarithmic differentiation with convergence controlled}

Fix a compact set disjoint from the negative zero set, and enclose it
in $|z|\leq r$. For all sufficiently large $n$,
$\gamma_n^2>2r$ and
\[
\left|\frac1{z+\gamma_n^2}\right|\leq\frac2{\gamma_n^2}.
\]
The right side is summable by \eqref{eq:rh-product}. Finite initial
factors have no zeros on the chosen compact set. Locally uniform
convergence of the product and its logarithmic derivatives therefore
justifies
\begin{equation}\label{eq:rh-log-derivative}
\frac{\Theta'(z)}{\Theta(z)}
=\sum_{n=1}^{\infty}\frac1{z+\gamma_n^2}
\end{equation}
away from the zeros. In particular this identity holds at every
$a\geq0$.

\begin{proposition}[Diagonal realization under RH]\label{prop:diagonal}
Let $\HH=\ell^2(\N;\C)$, and define
\begin{equation}\label{eq:diagonal-data}
(Kx)_n=\gamma_n^{-2}x_n,\qquad v_n=\gamma_n^{-1}.
\end{equation}
Then $K$ is bounded, self-adjoint, nonnegative, compact and trace
class; $v\in\HH$; and
\begin{equation}\label{eq:diagonal-calibration}
\ip v{(I+aK)^{-1}v}=L(a)\qquad(a\geq0).
\end{equation}
\end{proposition}
\begin{proof}
The positive ordinates have a positive minimum: a sequence tending to
zero would give an accumulation of zeros at $1/2$, contradicting
$\xi(1/2)>0$. Hence the diagonal entries of $K$ are bounded.
They are positive real numbers and tend to zero, which gives
self-adjointness, positivity and compactness. Their sum is finite
by \eqref{eq:rh-product}, so the positive diagonal operator is trace
class. The same sum equals $\sum_n|v_n|^2$, proving $v\in\HH$.
Finally,
\[
\ip v{(I+aK)^{-1}v}
=\sum_n\frac{\gamma_n^{-2}}{1+a\gamma_n^{-2}}
=\sum_n\frac1{a+\gamma_n^2}=L(a),
\]
using \eqref{eq:rh-log-derivative}.
\end{proof}

This proves $\RH\Rightarrow\HSeed$, completing
Theorem~\ref{thm:main}. Taking $a=0$ gives the normalization
\begin{equation}\label{eq:trace-norm}
\Tr K=\norm v^2=\sum_n\gamma_n^{-2}
=L(0)=\frac{\xi''(1/2)}{2\xi(1/2)}.
\end{equation}
For an arbitrary calibrated representation, $\norm v^2=L(0)$ still
holds. The equality with $\Tr K$ is asserted only for the particular
diagonal construction just given; arbitrary invisible operator components
need not have that trace.

If one prefers distinct positive ordinates $\gamma_j$ with multiplicities
$m_j$, the equivalent scalar model uses
$K_{jj}=\gamma_j^{-2}$ and $v_j=\sqrt{m_j}/\gamma_j$.
Neither version assumes that zeta zeros are simple. In the repeated
model, the vector need not be cyclic when an eigenvalue is repeated.
No cyclicity assumption is needed for the argument.

\section[Spectral measures and moments]{Spectral measures, moments and exact reconstruction}\label{sec:moments}

\subsection{A scalar measure formulation}

Define $\MSeed$ to mean that for some $M>0$ and finite positive
Borel measure $\mu$ on $[0,M]$,
\begin{equation}\label{eq:mseed}
L(a)=\int_{[0,M]}\frac1{1+at}\dd\mu(t)
\end{equation}
on a nonempty positive interval.
Section~\ref{sec:analytic} proves $\HSeed\Rightarrow\MSeed$.
Conversely, given \eqref{eq:mseed}, take
\[
\HH=L^2([0,M],\mu),\qquad(Ku)(t)=tu(t),\qquad v(t)=1.
\]
Finite mass makes $v$ square integrable, while multiplication by $t$
is bounded, nonnegative and self-adjoint. Its scalar resolvent is
exactly \eqref{eq:mseed}. Therefore
\begin{equation}\label{eq:measure-iff}
\MSeed\quad\Longleftrightarrow\quad\HSeed
\quad\Longleftrightarrow\quad\RH.
\end{equation}

The word ``finite'' in \emph{finite measure} means finite total mass.
It does not mean finite support, a finite list of rational numbers, or
a finite-checker witness. Under RH the measure is
\begin{equation}\label{eq:rh-measure}
\mu=\sum_n\gamma_n^{-2}\delta_{\gamma_n^{-2}},\qquad
\mu([0,M])=\kappa.
\end{equation}
This differs from the zero-counting measure
$\sum_n\delta_{\gamma_n^2}$, which has infinite total mass. The
reciprocal variable and the weights are essential to the normalization.

\subsection{Moments and positive finite sections}

For any fixed positive operator model, put
\begin{equation}\label{eq:moments}
m_n=\ip v{K^nv}=\int_{[0,M]}t^n\dd\mu(t)\qquad(n\geq0).
\end{equation}
If $|z|M<1$, the Neumann series converges uniformly in the spectral
variable and gives
\begin{equation}\label{eq:moment-series}
F(z)=\sum_{n=0}^{\infty}(-1)^nm_nz^n,
\qquad0\leq m_n\leq M^n m_0.
\end{equation}
The remainder after degree $N$ has the exact form
\begin{equation}\label{eq:moment-remainder}
F(z)-\sum_{n=0}^N(-1)^nm_nz^n
=(-z)^{N+1}\int\frac{t^{N+1}}{1+zt}\dd\mu(t),
\end{equation}
and consequently
\begin{equation}\label{eq:moment-bound}
\left|F(z)-\sum_{n=0}^N(-1)^nm_nz^n\right|
\leq\frac{m_0(M|z|)^{N+1}}{1-M|z|}.
\end{equation}
For a calibrated model,
\begin{equation}\label{eq:actual-moments}
m_n=\frac{(-1)^n}{n!}L^{(n)}(0).
\end{equation}
This is an all-orders identity for actual derivatives.

For a polynomial $p(t)=\sum_i c_it^i$, positivity of $\mu$ gives
\begin{align}
\sum_{i,j}\overline{c_i}c_jm_{i+j}&\geq0,\label{eq:H0}\\
\sum_{i,j}\overline{c_i}c_jm_{i+j+1}&\geq0,\label{eq:H1}\\
\sum_{i,j}\overline{c_i}c_j
\bigl(M^2m_{i+j}-m_{i+j+2}\bigr)&\geq0.\label{eq:H2}
\end{align}
These are the integrals of $|p|^2$, $t|p|^2$ and
$(M^2-t^2)|p|^2$. Every finite matrix in this hierarchy is positive
semidefinite. The following converse records exactly what an
\emph{all-orders} moment construction can accomplish.

\begin{theorem}[Bounded positive moment reconstruction]\label{thm:gns}
Let $(m_n)_{n\geq0}$ be a real sequence and $M>0$.
Suppose \eqref{eq:H0}--\eqref{eq:H2} hold for every finite coefficient
list. Then there exist a Hilbert space, a bounded nonnegative
self-adjoint operator $K$ with $\norm K\leq M$, and a vector $v$
such that $m_n=\ip v{K^nv}$ for all $n$.
\end{theorem}
\begin{proof}
Define a complex linear functional $\ell$ on polynomials by
$\ell(t^n)=m_n$, and a sesquilinear form by
$\ip pq_0=\ell(\overline p\,q)$, where conjugation acts on coefficients
and $t$ is real. Condition~\eqref{eq:H0} makes this form positive
semidefinite, hence it satisfies Cauchy--Schwarz. Quotient the polynomial
space by the vectors of zero seminorm and complete to a Hilbert space
$\HH$.

Multiplication by $t$ is well defined on the quotient: by
\eqref{eq:H2},
\[
\norm{tp}_0^2=\ell(t^2|p|^2)\leq M^2\ell(|p|^2)=M^2\norm p_0^2.
\]
Thus it sends every null vector to a null vector and is bounded by
$M$. It extends to an everywhere defined bounded operator $K$ on
$\HH$. Since
$\ell(\overline{tp}\,q)=\ell(\overline p\,tq)$, it is symmetric
on a dense subspace and hence self-adjoint after extension.
Condition~\eqref{eq:H1} gives
$\ip{[p]}{K[p]}=\ell(t|p|^2)\geq0$; continuity makes $K$ positive
on all of $\HH$. Let $v=[1]$. Then $K^nv=[t^n]$ and
$\ip v{K^nv}=\ell(t^n)=m_n$, as required.
\end{proof}

\begin{corollary}[An all-order actual-jet criterion]\label{cor:moment-criterion}
Define $m_n$ by \eqref{eq:actual-moments}. RH is equivalent to the
existence of $M>0$ such that \eqref{eq:H0}--\eqref{eq:H2} hold
for every finite coefficient list.
\end{corollary}
\begin{proof}
Under RH, the diagonal seed supplies these inequalities.
Conversely, Theorem~\ref{thm:gns} constructs $K,v$ with the specified
moments. The resolvent has Taylor series \eqref{eq:moment-series}.
The actual $L$ is analytic near zero and has exactly the same Taylor
coefficients. They therefore agree on a common disk of positive radius,
and in particular on a positive open interval. This is $\HSeed$,
which implies RH.
\end{proof}

The quantifier ``for every finite coefficient list'' ranges over
unbounded degrees. Checking finitely many matrices does not establish
this corollary. A finite certificate would have to derive the entire
hierarchy by an independently proved uniform mechanism.

\subsection{What is determined by calibration}

If two compactly supported finite positive measures represent the same
actual $L$ on an interval, analytic continuation gives equality near
zero. Their moments consequently agree. View both on a common compact
interval containing their supports. Polynomials are dense in the
continuous functions on that interval, and equality of integrals for
polynomials extends, by the total-mass bound, to all continuous functions.
Uniqueness of finite Borel measures then gives equality of the two measures.
Thus the scalar spectral measure seen by $v$ is determined. This does
not identify every operator on the ambient Hilbert space: an orthogonal
summand invisible to $v$ leaves all the scalar identities unchanged.

\section[Shifts and truncations]{Basepoint shifts and finite truncations}\label{sec:shifts}

\subsection{Two different shifted vectors}

Fix $a_0\geq0$, and abbreviate $R_0=(I+a_0K)^{-1}$ in this
subsection only. Set
\begin{equation}\label{eq:shift-data}
B=KR_0,\qquad w=R_0^{1/2}v,\qquad u=R_0v.
\end{equation}
These two vectors have different purposes. For $t=a-a_0$ with
$a\geq0$,
\[
I+aK=(I+a_0K)(I+tB).
\]
All the displayed operators commute. Therefore
\begin{equation}\label{eq:f-shift}
f(a_0+t)=\ip w{(I+tB)^{-1}w}.
\end{equation}
The functional calculus gives
$0\leq B\leq M(1+a_0M)^{-1}I$ whenever $0\leq K\leq MI$.
In particular the local shifted inverse is analytic at $t=0$.

On the other hand, \eqref{eq:second-resolvent} gives
\begin{equation}\label{eq:g-shift}
g(a_0+t)-g(a_0)
=t\ip u{(I+tB)^{-1}u}.
\end{equation}
To verify this, write
$aR_a-a_0R_0=tR_aR_0$ and
$R_a=(I+tB)^{-1}R_0$, then take the scalar product with $v$.
Differentiating \eqref{eq:g-shift} gives
\begin{equation}\label{eq:shift-derivative}
g'(a_0+t)=\norm{(I+tB)^{-1}u}^2.
\end{equation}
The additive value $g(a_0)$ and the distinction $w\ne u$ in general
must be retained. Merely replacing $a$ by $a-a_0$ in part of a
formula changes the mathematical condition.

\subsection{What finite diagonal truncations do prove}

Under RH, use the repeated-ordinate model from
Proposition~\ref{prop:diagonal}, and truncate after $N$ entries:
\[
F_N(z)=\sum_{n=1}^N\frac1{z+\gamma_n^2},\qquad
S_N=\sum_{n>N}\gamma_n^{-2}.
\]
For every $a\geq0$,
\begin{equation}\label{eq:positive-tail}
0\leq L(a)-F_N(a)\leq S_N.
\end{equation}
To state a bound on a whole disk, define the tail directly by
\[
T_N(z)=\sum_{n>N}\frac1{z+\gamma_n^2}.
\]
If $\gamma_{N+1}^2>2r$, this series is holomorphic on a neighborhood
of $|z|\leq r$. Away from the initial poles it equals $F-F_N$;
at a shared pole among the first $N$ terms, it is the holomorphic
extension of that difference after cancellation. Thus, on the entire disk,
\begin{equation}\label{eq:complex-tail}
|T_N(z)|\leq2S_N.
\end{equation}
For each derivative order $j\geq0$, the same argument gives
\begin{equation}\label{eq:derivative-tail}
|T_N^{(j)}(z)|
\leq j!\,2^{j+1}\sum_{n>N}\gamma_n^{-2j-2}.
\end{equation}
These estimates follow term by term from the convergent series and
its displayed denominator bounds. They show what a tail certificate
would need to control, including its dependence on derivative order.

The construction here assumes RH to identify the positive spectrum.
It is consequently not an unconditional finite computation proving RH.
Even a uniform small error does not imply exact equality; a uniform
error bound for values alone does not automatically supply an error
bound for all derivatives on the same domain.

\section[Hyperbolic zeta connections]{Hyperbolic zeta identities and the actual-seed transfer}\label{sec:hyperbolic}

The hyperbolic zeta identities provide explicit positive measures from
which the operator calculus of Section~\ref{sec:resolvent} can be
constructed. This section makes that connection precise, including the
normalization of the odd-zeta formulas. It also identifies a limitation:
the resulting model functions have a different large-parameter behavior
from the actual $L$. Their positivity is useful structural information,
but these particular models cannot themselves be its exact calibration.
We then derive an actual-$\xi$ Mellin identity and a conditional
coefficient-limit theorem for positive determinant models obtained
from the project's signed pencils. The limit condition is stated as
an open analytic obligation.

The source identities below are traced to the manuscript
\cite{hyp-manuscript} and repository revision \texttt{b56c810}
\cite{hyp-weight,hyp-moments,hyp-pole,hyp-gram,hyp-signs}.
Only the displayed identities are used, with proofs or explicit
classical inputs supplied here. Source inspection is separate from
native compilation and axiom auditing.

\subsection{A positive weight and its Mellin moments}

For an integer $k\geq0$, set
\begin{equation}\label{hyp-weight-def}
B(x)=\frac1{e^x-1},\qquad
\omega_k(y)=\frac{2(2\pi)^{2k}}{(2k)!}
y^{2k+1}B^{(2k)}(2\pi y),\qquad y>0.
\end{equation}
Here $B^{(2k)}$ denotes differentiation with respect to the argument of
$B$, not with respect to $y$. This distinction determines the factor
$(2\pi)^{2k}$. The weights and their moment identities are the relevant
positive-measure objects in the hyperbolic zeta development
\cite{hyp-weight,hyp-moments}. If a continuous weight parameter
$s>1$ is used, the coefficient extending these weights is
$2(2\pi)^{s-1}/\Gamma(s)$. The printed continuous coefficient
$(2\pi)^s/(2\Gamma(s))$ differs by $\pi/2$; the repository records
this correction explicitly \cite{hyp-normalization}.

\begin{proposition}[Positivity and exact moments]\label{hyp-moment-prop}
The function $\omega_k$ is strictly positive, tends to $1/\pi$ at zero,
and defines a finite measure on $(0,\infty)$. For every real $q>-1$,
\begin{equation}\label{hyp-mellin}
\int_0^\infty y^q\omega_k(y)\dd y
=\frac{2\Gamma(q+2k+2)\zeta(q+2)}
{(2k)!(2\pi)^{q+2}}.
\end{equation}
In particular,
\begin{equation}\label{hyp-mass}
\int_0^\infty\omega_k(y)\dd y=\frac{2k+1}{12}.
\end{equation}
\end{proposition}
\begin{proof}
The geometric series and its locally uniformly convergent derivatives
give
\[
B^{(2k)}(x)=\sum_{n=1}^{\infty}n^{2k}e^{-nx}>0\qquad(x>0).
\]
The Laurent expansion $B(x)=x^{-1}-1/2+O(x)$, with its analytic
remainder near zero, shows
$B^{(2k)}(x)\sim(2k)!x^{-2k-1}$. Substitution in
\eqref{hyp-weight-def} gives $\omega_k(y)\to1/\pi$.
At infinity the differentiated geometric series is $O(e^{-x})$,
so the weight has exponential decay times a polynomial.

For the exact value, Tonelli's theorem applies because every summand is
nonnegative. The gamma integral then yields
\begin{align*}
\int_0^\infty y^q\omega_k(y)\dd y
&=\frac{2(2\pi)^{2k}}{(2k)!}
\sum_{n\geq1}n^{2k}
\frac{\Gamma(q+2k+2)}{(2\pi n)^{q+2k+2}}\\
&=\frac{2\Gamma(q+2k+2)}{(2k)!(2\pi)^{q+2}}
\sum_{n\geq1}n^{-q-2}.
\end{align*}
The last series converges precisely in the stated real range. Taking
$q=0$ and using $\zeta(2)=\pi^2/6$ proves \eqref{hyp-mass}.
\end{proof}

Thus polynomial Gram matrices have exact entries involving even and odd
zeta values. Their positivity alone does not make multiplication by $y$
or $y^2$ bounded: the measure has unbounded support. The next construction
supplies the bounded operator required by the seed framework.
For instance, the even moments $b_j=\int y^{2j}\omega_k(y)\dd y$
satisfy
\begin{equation}\label{hyp-raw-ratio}
\frac{b_{j+1}}{b_j}
=\frac{(2k+2j+3)(2k+2j+2)}{4\pi^2}
\frac{\zeta(2j+4)}{\zeta(2j+2)}\longrightarrow\infty.
\end{equation}
Thus no fixed finite support bound can be read into the raw moment
sequence. Positivity and bounded support are different requirements in
the reconstruction theorem. Here the support is made compact by an
explicit change of spectral variable, whose effect on the vector and
measure must also be included.

\subsection{Compactification and the fixed Hilbert vector}

Put
\begin{equation}\label{hyp-compactification}
t(y)=\frac1{1+y^2},\qquad
\HH_k=L^2((0,\infty),\omega_k(y)\dd y),\qquad
(K_ku)(y)=t(y)u(y),\qquad v_k(y)=\sqrt{t(y)}.
\end{equation}
Then $K_k$ is nonnegative and self-adjoint, with $\norm{K_k}=1$.
Indeed, $0<t(y)<1$, and values arbitrarily close to $1$ occur on
sets of positive measure. Also $v_k\in\HH_k$ by
\eqref{hyp-mass}. Its scalar resolvent is
\begin{equation}\label{hyp-resolvent}
F_k(a)=\ip{v_k}{(I+aK_k)^{-1}v_k}
=\int_0^\infty\frac{\omega_k(y)}{a+1+y^2}\dd y,
\qquad a>-1.
\end{equation}
In particular, $\norm{v_k}^2=F_k(0)$; it is not the total mass in
\eqref{hyp-mass}.
The factor $\sqrt{t(y)}$ in the vector is essential to this
normalization: its square supplies the numerator $t(y)$ in
$t(y)/(1+at(y))=(a+1+y^2)^{-1}$. All parameters consequently use
one operator and one vector, exactly as required in the seed definitions.
The construction is infinite dimensional, despite the single integer
parameter specifying the density.

Define the finite measure $\nu_k$ on $[0,1]$ by
\[
\int\varphi(t)\dd\nu_k(t)
=\int_0^\infty t(y)\varphi(t(y))\omega_k(y)\dd y.
\]
Its moments are
\begin{equation}\label{hyp-compact-moments}
q_n^{(k)}=\int_0^1t^n\dd\nu_k(t)
=\int_0^\infty\frac{\omega_k(y)}{(1+y^2)^{n+1}}\dd y,
\qquad
F_k(a)=\sum_{n\geq0}(-1)^nq_n^{(k)}a^n\quad(|a|<1).
\end{equation}
The series follows from a uniformly convergent geometric series on
$0\leq t\leq1$. For every polynomial $p(t)=\sum_i c_it^i$,
integration of $|p|^2$, $t|p|^2$, $(1-t)|p|^2$ and
$(1-t^2)|p|^2$ proves positivity of all finite sections of
\begin{equation}\label{hyp-psd}
(q_{i+j}^{(k)}),\quad(q_{i+j+1}^{(k)}),\quad
(q_{i+j}^{(k)}-q_{i+j+1}^{(k)}),\quad
(q_{i+j}^{(k)}-q_{i+j+2}^{(k)}).
\end{equation}
These are all-order identities for this explicitly defined model.
They make no identification with the actual moments
$(-1)^nL^{(n)}(0)/n!$.
There is also the full compact-moment difference identity
\begin{equation}\label{hyp-hausdorff}
\sum_{j=0}^{r}(-1)^j\binom rj q_{n+j}^{(k)}
=\int_0^1 t^n(1-t)^r\dd\nu_k(t)\geq0
\qquad(n,r\geq0).
\end{equation}
For a nonzero polynomial, each quadratic form in
\eqref{hyp-psd} is strictly positive: its integrand is positive
except at finitely many roots and possibly the endpoints, and the
transformed measure gives positive mass to every open subinterval
of $(0,1)$.

The compact moments should not be obtained by formally expanding
$(1+y^2)^{-1}$ into powers of $y^2$ and integrating term by term.
That geometric series does not converge throughout the original
unbounded support. Formula~\eqref{hyp-compact-moments}, or the
Hurwitz calculation below, supplies a justified route instead.

With $R_{k,a}=(I+aK_k)^{-1}$ and $g_k(a)=aF_k(a)$,
the general resolvent calculation gives the explicit formulas
\begin{align}
g_k'(a)
&=\norm{R_{k,a}v_k}^2
=\int_0^\infty\frac{(1+y^2)\omega_k(y)}{(a+1+y^2)^2}\dd y,
\label{hyp-energy}\\
\frac{g_k(a)-g_k(b)}{a-b}
&=\ip{R_{k,a}v_k}{R_{k,b}v_k}
=\int_0^\infty
\frac{(1+y^2)\omega_k(y)}{(a+1+y^2)(b+1+y^2)}\dd y.
\label{hyp-gram}
\end{align}
The second formula is for distinct nonnegative $a,b$ and extends
diagonally by the first. Multiplication by $4\sqrt{ab}$ produces
the Gram identity for $2\sqrt a\,R_{k,a}v_k$.

\subsection{The Hurwitz identity and its endpoint normalization}

\begin{proposition}[Real-parameter odd-zeta resolvent]\label{hyp-hurwitz-prop}
For $k\geq1$, $a>-1$, and $z=\sqrt{1+a}>0$,
\begin{equation}\label{hyp-hurwitz}
F_k(a)=z^{2k}\zeta(2k+1,z)-\frac1{2z}-\frac1{2k}.
\end{equation}
For $k=0$,
\begin{equation}\label{hyp-digamma}
F_0(a)=\log z-\frac1{2z}-\psi(z).
\end{equation}
\end{proposition}
\begin{proof}
Binet's digamma formula \cite{dlmf-binet} states that, for real $z>0$,
\begin{equation}\label{hyp-binet-base}
h(z):=\log z-\frac1{2z}-\psi(z)
=2\int_0^\infty\frac{yB(2\pi y)}{y^2+z^2}\dd y.
\end{equation}
This immediately proves the case $k=0$. For $k\geq1$, use
$\partial_y^{2k}B(2\pi y)=(2\pi)^{2k}B^{(2k)}(2\pi y)$
and integrate by parts $2k$ times:
\[
F_k(z^2-1)=\frac2{(2k)!}\int_0^\infty B(2\pi y)
\partial_y^{2k}\!\left(\frac{y^{2k+1}}{y^2+z^2}\right)\dd y.
\]
Every boundary term vanishes. At zero, the product of the relevant
derivatives is $O(y)$: the rational factor has a zero of order
$2k+1$, whereas $\partial_y^jB(2\pi y)=O(y^{-j-1})$.
At infinity the exponential decay dominates the polynomial factors.

Polynomial division gives
\[
\frac{y^{2k+1}}{y^2+z^2}
=P_{2k-1}(y)+(-z^2)^k\frac{y}{y^2+z^2}.
\]
The polynomial disappears after $2k$ differentiations. From
$y/(y^2+z^2)=\tfrac12((y+iz)^{-1}+(y-iz)^{-1})$ one obtains
\[
\partial_y^{2k}\frac{y}{y^2+z^2}
=(-1)^k\partial_z^{2k}\frac{y}{y^2+z^2}.
\]
On compact positive $z$-intervals these integral derivatives are dominated
near zero and at infinity. Consequently
\begin{equation}\label{hyp-binet-differentiated}
F_k(z^2-1)=\frac{z^{2k}}{(2k)!}h^{(2k)}(z).
\end{equation}
More explicitly, every fixed $z$-derivative of
$y/(y^2+z^2)$ is $O(y)$ near $y=0$, uniformly when $z$ stays
away from zero. Multiplication by $B(2\pi y)=O(y^{-1})$
therefore gives a locally integrable uniform bound. At infinity,
the Bose factor supplies exponential domination. This verifies the
exchange of differentiation and integration used in
\eqref{hyp-binet-differentiated}.
Finally,
$\psi^{(2k)}(z)=-(2k)!\zeta(2k+1,z)$
\cite{dlmf-hurwitz}; differentiating the logarithm and reciprocal in
$h$ supplies respectively $-1/(2k)$ and $-1/(2z)$ after multiplication
by $z^{2k}/(2k)!$. This proves \eqref{hyp-hurwitz}.
\end{proof}

Taking values and first derivatives at $a=0$ gives
\begin{align}
q_0^{(k)}&=\zeta(2k+1)-\frac12-\frac1{2k},\notag\\
q_1^{(k)}&=\frac{2k+1}{2}\zeta(2k+2)
-k\zeta(2k+1)-\frac14\qquad(k\geq1).
\label{hyp-first-moments}
\end{align}
For $k=0$ the corresponding values are
$q_0^{(0)}=\gamma-1/2$ and
$q_1^{(0)}=\zeta(2)/2-3/4$, where $\gamma$ is Euler's constant.
The differentiation uses $\partial_a=(2z)^{-1}\partial_z$ and
$\partial_z\zeta(s,z)=-s\zeta(s+1,z)$. For $k\geq1$, repetition
expresses every fixed compact moment through finitely many adjacent zeta values and
rational constants. The integral formulas guarantee, independently of
any cancellation in those expressions, that
$q_0^{(k)}>q_1^{(k)}>q_2^{(k)}>\cdots>0$.
This gives exact positivity identities suitable for comparing proposed
moment constructions, without supplying the actual $\Theta$ jets.

The integer-parameter versions must retain their endpoint convention
\cite{hyp-pole}. For $k\geq1$, $N\geq1$, and $s=2k+1$,
\begin{align}
F_k(N^2-1)
&=N^{2k}\left(\zeta(s)-\sum_{n=1}^{N-1}n^{-s}\right)
-\frac1{2N}-\frac1{2k}\notag\\
&=N^{2k}\left(\zeta(s)-\sum_{n=1}^{N}n^{-s}\right)
+\frac1{2N}-\frac1{2k}.
\label{hyp-integer-endpoint}
\end{align}
Including the $N$th summand changes the sign of the displayed
$1/(2N)$ contribution. At $N=1$, either version recovers
the value of $q_0^{(k)}$ in \eqref{hyp-first-moments}.

\subsection{An exact Bose term in the actual logarithmic derivative}

There is a direct identity involving the actual completed zeta function.
Let $a>1/4$ and $s=1/2+\sqrt a$, so $s>1$. Logarithmic
differentiation of \eqref{eq:xi}, Binet's formula, and the convergent
von Mangoldt series \cite{dlmf-binet,dlmf-mangoldt} give
\begin{align}
2\sqrt a\,L(a)
={}&\frac1{2s}+\frac1{s-1}+\frac12\log\frac{s}{2\pi}\notag\\
&-\int_0^\infty
\frac{y\dd y}{(y^2+s^2/4)(e^{2\pi y}-1)}
-\sum_{n=2}^{\infty}\Lambda(n)n^{-s}.
\label{hyp-actual-binet}
\end{align}
Indeed, before substitution the gamma contribution is
$\tfrac12\psi(s/2)$ and the elementary contribution includes $1/s$.
The complete starting identity is
\[
\frac{\xi'(s)}{\xi(s)}
=\frac1s+\frac1{s-1}-\frac12\log\pi
+\frac12\psi(s/2)+\frac{\zeta'(s)}{\zeta(s)}.
\]
Binet contributes $-1/(2s)$, leaving the coefficient $1/(2s)$
in \eqref{hyp-actual-binet}. The same Bose kernel therefore occurs
in the actual function, but together with the logarithmic, rational,
and prime contributions. The positivity of one integral does not
establish a positive-resolvent representation of their complete sum.

\begin{proposition}[Asymptotic obstruction to direct calibration]
\label{hyp-obstruction}
No $F_k$ equals the actual $L$ on a nonempty open positive interval.
More generally, this remains impossible for a finite sum
$G(a)=\sum_{j=1}^{J}c_jF_{k_j}(d_ja+e_j)$ with
$c_j,d_j>0$ and $e_j\geq0$.
\end{proposition}
\begin{proof}
Dominated convergence in \eqref{hyp-resolvent} gives
\begin{equation}\label{hyp-model-asymptotic}
\lim_{a\to\infty}aF_k(a)=\frac{2k+1}{12},\qquad
\lim_{a\to\infty}aG(a)=
\sum_{j=1}^{J}\frac{c_j(2k_j+1)}{12d_j}<\infty.
\end{equation}
For the actual function, the integral in \eqref{hyp-actual-binet}
is bounded by $1/(6s^2)$, using
$\int_0^\infty yB(2\pi y)\dd y=1/24$.
The prime sum is exponentially small as $s\to\infty$:
for $s\geq4$ it is bounded by
$2^{-s/2}\sum_{n\geq2}(\log n)n^{-2}$.
The remaining rational terms are $O(s^{-1})$. Since
$s=1/2+\sqrt a$, it follows that
\begin{equation}\label{hyp-actual-asymptotic}
L(a)=\frac{\log a}{8\sqrt a}
-\frac{\log(2\pi)}{4\sqrt a}+O(a^{-1}),
\qquad aL(a)\longrightarrow+\infty.
\end{equation}
Both $G$ and $L$ are real analytic throughout $(0,\infty)$.
Equality on an open subinterval would imply equality on that connected
ray, contradicting \eqref{hyp-model-asymptotic} and
\eqref{hyp-actual-asymptotic}.
\end{proof}

This obstruction is consistent with the abstract Hilbert-seed theorem.
A finite compact spectral measure need not have a finite reciprocal
moment. For the present models, however,
\begin{equation}\label{hyp-reciprocal-moment}
\int_{(0,1]}\frac1t\dd\nu_k(t)
=\int_0^\infty\omega_k(y)\dd y=\frac{2k+1}{12}.
\end{equation}
This extra integrability is exactly what makes $aF_k(a)$ tend to
a finite limit. Under RH, the canonical measure in
\eqref{eq:rh-measure} instead has reciprocal moment
$\sum_n1=\infty$, although its total mass is finite. Thus one
cannot reject all bounded positive operators from
Proposition~\ref{hyp-obstruction}. The result rejects this particular
specific Bose multiplication construction and its stated finite combinations.
It identifies a concrete property that a different construction would
have to change.

Similarly, the positive matrices in \eqref{hyp-psd} can serve as exact
examples for checking normalization, shifts and proposed reconstruction
lemmas. Their proofs work uniformly in degree because a specific positive
measure was supplied first. Transferring that conclusion to the actual
jets would require a theorem equating the entire moment sequences, not
agreement of several low-order entries or a numerical fit on an interval.

The usable conclusion is precise: the hyperbolic identities supply
explicit positive models, exact moment hierarchies, and a Bose integral
inside the actual calibration problem. The bounded compactification
above does not solve that problem. Any further construction must account
for the full expression \eqref{hyp-actual-binet}, including the growth
in \eqref{hyp-actual-asymptotic}, before it can satisfy $\HSeed$.


\subsection{From hyperbolic source identities to the actual logarithmic moments}
\label{hyp-sec:actual-moments}

The useful connection begins with an identity for the actual completed
function, rather than with a positive model chosen independently of it.
Write $B(x)=(e^x-1)^{-1}$ for $x>0$. Elementary cancellation gives the
positive density
\begin{equation}\label{hyp-eq:fermi-density}
h(x):=4B(2x)-8B(4x)
=\frac{4}{e^{2x}+1}=2(1-\tanh x)>0.
\end{equation}
Thus a difference of two Bose kernels is exactly a positive Fermi,
or hyperbolic-tangent, kernel. Its positivity follows after the
cancellation; the two separate terms should not be treated as a
positive-measure decomposition.

Define the Mellin integral
\begin{equation}\label{hyp-eq:mellin}
I(s)=\int_0^\infty x^{s-1}h(x)\,\mathrm dx
=2^{2-s}\Gamma(s)(1-2^{1-s})\zeta(s),
\qquad \Re s>0.
\end{equation}
Here the expression at $s=1$ is removable, with $I(1)=2\log2$.
For $\Re s>1$, expand
$(e^{2x}+1)^{-1}=\sum_{k\geq1}(-1)^{k-1}e^{-2kx}$.
Absolute convergence after integration permits termwise use of the
gamma integral and gives
$I(s)=2^{2-s}\Gamma(s)\sum_{k\geq1}(-1)^{k-1}k^{-s}$.
Separating even and odd terms proves \eqref{hyp-eq:mellin} there.
Both sides continue holomorphically to $\Re s>0$: on the integral
side, $h(x)=O(1)$ at zero and $h(x)=O(e^{-2x})$ at infinity;
on the function side, the factor $1-2^{1-s}$ removes the pole of
$\zeta$ at $1$. The identity theorem gives the claimed domain.
This is the scaled Fermi integral in \cite[Eq.~25.5.3]{dlmf-fermi}.

The same density is obtained from the explicitly defined two-variable measure
\begin{equation}\label{hyp-eq:gp-pushforward}
q(r)=\frac{r}{e^r-1},\qquad d(r)=-q'(r)>0,\qquad
\mathrm d\mu(r,y)=2d(r)\,\mathrm dr\,\frac{\mathrm dy}{y},
\quad \frac14\leq y\leq\frac12,
\end{equation}
by pushing forward under $x=ry$. Indeed its density is
\[
2\int_{1/4}^{1/2}d(x/y)\,\frac{\mathrm dy}{y^2}
=\frac2x\bigl(q(2x)-q(4x)\bigr)=h(x).
\]
For completeness, $d(r)>0$ because
$-q'(r)=[e^r(r-1)+1]/(e^r-1)^2$, and the numerator vanishes at
zero and has derivative $re^r>0$. This calculation gives an exact
one-variable realization of the source layer. In particular its total
mass is $I(1)=2\log2$, and the original two-variable integral of
$(ry)^{s-1}$ is precisely $I(s)$. Tonelli's theorem applies first to
nonnegative integrands; absolute integrability then gives the same
change of variables for the complex Mellin integrands used here.

The gamma duplication formula \cite{dlmf-duplication} simplifies the
normalizing factor to
\begin{equation}\label{hyp-eq:A}
A(s)=\frac{s(s-1)\pi^{(1-s)/2}}
{4\Gamma((s+1)/2)(1-2^{1-s})},
\qquad \xi(s)=A(s)I(s).
\end{equation}
The equality holds in a complex neighborhood of $s=1/2$, which is
the domain needed here. The denominator is nonzero there, and
$A$ is positive on a real neighborhood of $1/2$: both $s(s-1)$
and $1-2^{1-s}$ are negative for $0<s<1$.
We do not claim that the displayed $A$ is holomorphic throughout
the right half-plane; its other apparent singularities require
cancellation in the product.

\subsubsection{An explicit map from the source integrals to the required moments}

Differentiation under the integral is valid to every fixed order:
\begin{equation}\label{hyp-eq:log-integrals}
I_j(s):=I^{(j)}(s)
=\int_0^\infty x^{s-1}(\log x)^j h(x)\,\mathrm dx.
\end{equation}
On each compact substrip $\Re s\geq\sigma>0$, the behavior at
zero is controlled by $x^{\sigma-1}|\log x|^j$, and the exponential
tail controls infinity. In particular all the integrals at $1/2$
are absolutely convergent. The product rule now gives actual,
unconditional coefficient identities:
\begin{align}
a_n&:=\frac{\xi^{(2n)}(1/2)}{(2n)!\,\xi(1/2)},
\qquad a_0=1,\label{hyp-eq:actual-coefficients}\\
\xi^{(2n)}(1/2)
&=\sum_{j=0}^{2n}\binom{2n}{j}
A^{(2n-j)}(1/2)I_j(1/2),\label{hyp-eq:jet-product}\\
\frac{\Theta(z)}{\Theta(0)}&=\sum_{n\geq0}a_nz^n.
\label{hyp-eq:normalized-theta-series}
\end{align}
Consequently, hyperbolic identities that evaluate or bound these
integrals can be transported to the actual $\Theta$ coefficients
with every normalization factor visible.

The resolvent moments are different quantities. Set
$m_n=(-1)^nL^{(n)}(0)/n!$ as before. Comparing coefficients in
$\Theta'=L\Theta$ gives the exact triangular recurrence
\begin{equation}\label{hyp-eq:newton-recurrence}
(n+1)a_{n+1}=\sum_{k=0}^{n}(-1)^km_k a_{n-k}
\qquad(n\geq0).
\end{equation}
It determines every required $m_n$ from $a_1,\ldots,a_{n+1}$.
The first cases and the first nontrivial Hankel determinant are
\begin{align}
m_0&=a_1,\notag\\
m_1&=a_1^2-2a_2,\notag\\
m_2&=a_1^3-3a_1a_2+3a_3,\label{hyp-eq:first-moments}\\
\det\begin{pmatrix}m_0&m_1\\m_1&m_2\end{pmatrix}
&=a_1^2a_2+3a_1a_3-4a_2^2.
\label{hyp-eq:first-hankel}
\end{align}
These are algebraic identities, not positivity assertions.
In particular, positive formulas for the individual $a_n$ do not
settle the signs of their nonlinear combinations. To check the
recurrence directly, multiply
$\sum_{k\geq0}(-1)^km_kz^k$ by
$\sum_{j\geq0}a_jz^j$ in a sufficiently small disk where
$\Theta$ is nonzero. The coefficient of $z^n$ is the right side
of \eqref{hyp-eq:newton-recurrence}, whereas differentiation of the
second series gives its left side. The occurrence of $m_n$ with
coefficient $(-1)^na_0=(-1)^n$ makes the recursion unique.

There is also a compact cumulant interpretation. Let $X=\log x$
under the probability measure
$\mathrm d\nu=x^{-1/2}h(x)\,\mathrm dx/I(1/2)$, and write
$\operatorname{cum}_j(X)$ for its $j$th cumulant. Its moment
generating function exists near zero, and
$\operatorname{cum}_j(X)=(\log I)^{(j)}(1/2)$.
Indeed $\mathbf E_\nu[e^{tX}]=I(1/2+t)/I(1/2)$ for
$|\Re t|<1/2$. The preceding endpoint estimates control this
identity and all its fixed derivatives on smaller strips.
Evenness of $\xi(1/2+w)$ and its nonzero value at zero give
\begin{equation}\label{hyp-eq:cumulants}
m_n=\frac{(-1)^n}{2(2n+1)!}
\left((\log A)^{(2n+2)}(1/2)
+\operatorname{cum}_{2n+2}(X)\right).
\end{equation}
All logarithms here are local analytic logarithms chosen real on
the relevant real interval. Formula~\eqref{hyp-eq:cumulants}
shows explicitly why ordinary positive-measure moments do not
automatically yield the required alternating logarithmic moments.
The derivatives of $\log A$ are part of the calibration and cannot
be discarded. For example, the second derivative of $\log I$ is
the variance of $X$, but the second derivative of $\log\xi$
also contains $(\log A)''$. Higher cumulants have no positivity
law comparable to the elementary squared-norm identity for ordinary
moment matrices.

\subsubsection{The positive mass and the remaining positivity theorem}

The mass $\kappa=m_0$ is strictly positive without RH. One standard
positive Riemann kernel, obtained by rescaling the representation in
\cite[Section~4, Eqs.~(4.1)--(4.2)]{csordas-kernel}, is
\begin{align}
\Psi(u)&=4\sum_{n\geq1}\pi n^2
\bigl(2\pi n^2e^{2u}-3\bigr)
e^{5u/2-\pi n^2e^{2u}},\qquad u\geq0,
\label{hyp-eq:positive-riemann-kernel}\\
\xi(1/2+w)&=\int_0^\infty\Psi(u)\cosh(wu)\,\mathrm du.
\label{hyp-eq:riemann-cosh}
\end{align}
Every summand is positive, since $2\pi n^2e^{2u}>3$.
The rapid decay justifies differentiation, and therefore
\begin{equation}\label{hyp-eq:kappa-positive}
\kappa=\frac{\int_0^\infty u^2\Psi(u)\,\mathrm du}
{2\int_0^\infty\Psi(u)\,\mathrm du}>0.
\end{equation}

To see the distinction quantitatively, normalize
$\mathrm d\nu_R=\Psi(u)\,\mathrm du/\xi(1/2)$ and put
$b_j=\int u^{2j}\,\mathrm d\nu_R$, so that
$a_j=b_j/(2j)!$ and $b_0=1$. Then
\begin{equation}\label{hyp-eq:kurtosis-gap}
m_1=\frac{3b_1^2-b_2}{12}.
\end{equation}
Ordinary positive-measure Cauchy--Schwarz gives $b_2\geq b_1^2$;
the necessary resolvent inequality $m_1\geq0$ instead requires
the additional upper bound $b_2\leq3b_1^2$. These are different
inequalities. The first Hankel determinant in
\eqref{hyp-eq:first-hankel} introduces still another combination
involving $b_3$, and the hierarchy continues to unbounded degree.

Such kernel positivity must still be distinguished from the full
resolvent criterion. For example, the entire function
$\Theta_*(z)=\tfrac34+\tfrac14\cosh\sqrt z$ has a positive
cosh-measure representation and is positive for $z\geq0$.
Nevertheless its normalized coefficients satisfy
$a_1=1/8$, $a_2=1/96$, and its transformed moment is
$m_1=a_1^2-2a_2=-1/192<0$.
Thus even a positive cosh representation and entire square descent
do not by themselves imply the desired Stieltjes moment positivity.

The exact benefit of the hyperbolic source identities is therefore
the calibrated chain
$I_j\longmapsto\xi^{(2n)}(1/2)\longmapsto a_n\longmapsto m_n$.
The missing theorem must control these \emph{actual transformed}
moments: for one finite $M>0$, all finite sections of
$(m_{i+j})$, $(m_{i+j+1})$, and
$(M^2m_{i+j}-m_{i+j+2})$ must be positive semidefinite.
The source integrals already give positivity of ordinary
log-moment matrices $(I_{i+j}(1/2))$, since for every finite
coefficient list
\[
\sum_{i,j}\overline c_i c_j I_{i+j}(1/2)
=\int_0^\infty x^{-1/2}h(x)
\left|\sum_jc_j(\log x)^j\right|^2\,\mathrm dx\geq0.
\]
These are moments of $\log x$, which ranges over the whole real
line. They are not already moments of a bounded nonnegative
operator with the required logarithmic-derivative calibration.
Transferring their positivity
through \eqref{hyp-eq:jet-product} and
\eqref{hyp-eq:newton-recurrence} is a substantive additional
obligation. Establishing it uniformly would activate the paper's
bounded moment reconstruction and RH transfer.


\subsection{From the signed hyperbolic pencil to positive determinant models}
\label{subsec:hyp-pencil}

The determinant identities in the Hyperbolic Zeta project give a second
construction, different from the multiplication model above. Let
$k,m,N,J$ be integers with $k,m\geq1$ and $1\leq N<J$, and put
$h=J-N$, $x_0=\zeta(2k+1)$ and
\[
D_j(t)=\prod_{r=1}^{j}(t+r^2).
\]
For an admissible affine Hankel pencil in that project, its value $G$
at $x_0$ has entries
\begin{equation}\label{eq:hyp-pencil-gram}
G_{ij}=\int_0^\infty y^{2(i+j)}
  \frac{D_N(y^2)^{2m}}{D_J(y^2)}\omega_k(y)\dd y,
\qquad 0\leq i,j<h.
\end{equation}
All these integrals converge. For $c\in\C^h\setminus\{0\}$,
\[
c^*Gc=\int_0^\infty
 \left|\sum_{j=0}^{h-1}c_j y^{2j}\right|^2
 \frac{D_N(y^2)^{2m}}{D_J(y^2)}\omega_k(y)\dd y>0.
\]
The density and rational factor are positive, and a nonzero polynomial
cannot vanish on an interval. Thus $G$ is positive definite.
The affine pencil is
\begin{equation}\label{eq:hyp-affine-pencil}
\mathcal G(X)=G+(X-x_0)M,\qquad
\Delta(X)=\det\mathcal G(X),\qquad
B=G^{-1/2}MG^{-1/2}.
\end{equation}
Here $M=M^*$ is the coefficient of $X$. These identifications are
recorded in the project's exact Gram and pole-diagonalization sources
\cite{hyp-gram,hyp-signs}.

The positive Gram matrix does not make the coefficient $M$ positive.
In the pole basis the identity is
\begin{equation}\label{eq:hyp-pole-sign}
M=V\operatorname{diag}(\gamma_r)_{N<r\leq J}V^T,
\qquad \operatorname{sgn}\gamma_r=(-1)^{r-1},
\end{equation}
with $V$ invertible. In fact,
\[
\gamma_r=r^{2k}
 \frac{D_N(-r^2)^{2m-1}}{D_{\mathrm{tail}}'(-r^2)},
\qquad D_{\mathrm{tail}}(t)=\prod_{j=N+1}^J(t+j^2).
\]
The numerator has sign $(-1)^N$ and the denominator has sign
$(-1)^{r-N-1}$, proving the sign in \eqref{eq:hyp-pole-sign}.
Consequently $B$ is self-adjoint and invertible, but has both positive
and negative eigenvalues when $h\geq2$. Hyperbolicity of $\Delta$
means its roots are real; the Hilbert seed instead requires a
nonnegative operator. These are different properties.

\begin{proposition}[Squared-pencil identity]\label{prop:hyp-square}
For any finite self-adjoint $B$ as in \eqref{eq:hyp-affine-pencil}
and any $c>0$, let $K=cB^2$ and $Q_c(z)=\det(I+zK)$. Then
\begin{equation}\label{eq:hyp-square-product}
Q_c(z)=
\frac{\Delta(x_0+i\sqrt{cz})\Delta(x_0-i\sqrt{cz})}
     {\Delta(x_0)^2}.
\end{equation}
The right side is a polynomial in $z$, independent of the choice
of square root. If $\beta_1,\ldots,\beta_h$ are the eigenvalues of $B$,
then $\lambda_j=c\beta_j^2\geq0$ and
\begin{equation}\label{eq:hyp-square-resolvent}
Q_c(z)=\prod_{j=1}^h(1+z\lambda_j),\qquad
\frac{Q_c'(a)}{Q_c(a)}
=\sum_{j=1}^h\frac{\lambda_j}{1+a\lambda_j}
=\ip{v}{(I+aK)^{-1}v},\quad a\geq0,
\end{equation}
where $v=\sum_j\sqrt{\lambda_j}\,e_j$ in an orthonormal eigenbasis
of $K$. In particular the exact differential seed algebra holds for
$Q_c'/Q_c$.
\end{proposition}

\begin{proof}
Congruence in \eqref{eq:hyp-affine-pencil} gives
$\Delta(x_0+u)/\Delta(x_0)=\det(I+uB)$.
The commuting factors satisfy
\[
(I+i\sqrt{cz}B)(I-i\sqrt{cz}B)=I+czB^2.
\]
Taking determinants proves \eqref{eq:hyp-square-product}. The spectral
theorem gives the product and the scalar resolvent in
\eqref{eq:hyp-square-resolvent}. Applying the already proved identity
\eqref{eq:norm-square} to this $K,v$ gives the differential identity.
For real $z\geq0$, the numerator in \eqref{eq:hyp-square-product} is
an absolute square; for general complex $z$ it is the product of the
two displayed factors.
\end{proof}

The operation $B\mapsto B^2$ changes the scalar determinant. It
therefore constructs positive candidates without transferring any
identity with the actual $\Theta$. The next result states a sufficient
transfer condition in terms of exactly specified coefficients.

\subsection{An all-order coefficient-limit theorem}
\label{subsec:hyp-limit}

Write $\Theta(z)/\Theta(0)=\sum_{r\geq0}a_rz^r$, where $a_0=1$,
and let $\kappa=a_1=L(0)>0$ as established above. Consider a sequence
of finite, nonzero self-adjoint matrices $B_n$, for example a
specified sequence of the hyperbolic slopes. Define
\begin{equation}\label{eq:hyp-normalized-square}
c_n=\frac{\kappa}{\Tr(B_n^2)},\qquad
K_n=c_nB_n^2,\qquad
Q_n(z)=\det(I+zK_n)=\sum_{r\geq0}e_{n,r}z^r.
\end{equation}
Coefficients above the degree of $Q_n$ are set equal to zero.
The denominator is positive, and
\begin{equation}\label{eq:hyp-trace-bound}
K_n\geq0,\qquad \Tr K_n=\kappa,\qquad
\norm{K_n}\leq\kappa,\qquad
0\leq e_{n,r}\leq\frac{\kappa^r}{r!}.
\end{equation}
The last inequality follows by expanding $(\sum_j\lambda_{n,j})^r$:
the terms with $r$ distinct indices include $r!$ copies of every
summand of the elementary symmetric polynomial $e_{n,r}$.

\begin{theorem}[Conditional transfer of hyperbolic determinant models]
\label{thm:hyp-coefficient-transfer}
For the sequence \eqref{eq:hyp-normalized-square}, the following two
conditions are equivalent:
\begin{align}
&\text{for every }r\geq0,\qquad
\lim_{n\to\infty}e_{n,r}
=\frac{\xi^{(2r)}(1/2)}{(2r)!\,\xi(1/2)}=a_r;
\label{eq:hyp-coefficient-target}\\
&\text{for every }r\geq0,\qquad
\lim_{n\to\infty}\Tr(K_n^{r+1})
=\frac{(-1)^r}{r!}L^{(r)}(0)=m_r.
\label{eq:hyp-power-target}
\end{align}
Either condition implies locally uniform convergence on $\C$,
\begin{equation}\label{eq:hyp-entire-limit}
Q_n(z)\longrightarrow\frac{\Theta(z)}{\Theta(0)},
\end{equation}
and implies RH. No assertion that the particular Hyperbolic Zeta
sequence satisfies either condition is included in this theorem.
\end{theorem}

\begin{proof}
Put $p_{n,j}=\Tr(K_n^j)$ for $j\geq1$. Newton's identities give
\[
r e_{n,r}=\sum_{j=1}^r(-1)^{j-1}e_{n,r-j}p_{n,j}.
\]
The actual coefficients satisfy the identical finite algebraic
recurrence
\[
r a_r=\sum_{j=1}^r(-1)^{j-1}a_{r-j}m_{j-1},
\]
obtained from $\Theta'=L\Theta$. For any fixed $r$, the first
recurrence determines $e_{n,r}$ from $p_{n,1},\ldots,p_{n,r}$ and
also determines $p_{n,r}$ from $e_{n,1},\ldots,e_{n,r}$, since
$e_{n,0}=1$. Induction on $r$ proves the equivalence of
\eqref{eq:hyp-coefficient-target} and \eqref{eq:hyp-power-target}.

Assume coefficient convergence. Taking limits in
\eqref{eq:hyp-trace-bound} gives $0\leq a_r\leq\kappa^r/r!$.
On every disk $|z|\leq R$, both coefficient series are dominated by
$\sum_r(\kappa R)^r/r!=e^{\kappa R}$. Split the series at a fixed
degree, use coefficient convergence for its finite head, and then
let that degree tend to infinity. This proves
\eqref{eq:hyp-entire-limit}. Quantitatively, for every integer $d\geq0$,
\begin{equation}\label{eq:hyp-uniform-tail}
\sup_{|z|\leq R}\left|Q_n(z)-\frac{\Theta(z)}{\Theta(0)}\right|
\leq\sum_{r=0}^{d}|e_{n,r}-a_r|R^r
 +\frac{2e^{\kappa R}(\kappa R)^{d+1}}{(d+1)!}.
\end{equation}

Every zero of $Q_n$ is $-1/\lambda_{n,j}$ for a positive eigenvalue.
Since $\lambda_{n,j}\leq\kappa$, all $Q_n$ are zero-free on the
connected domain $\C\setminus(-\infty,-1/\kappa]$. Their locally
uniform limit equals $1$ at the origin and is not identically zero.
Hurwitz's theorem therefore makes this limit zero-free on the same
domain. The zeros of $\Theta$ lie on the negative real axis, giving RH
by Section~\ref{sec:theta}.

There is also a direct route through the moment criterion. The measures
$\mu_n=\sum_j\lambda_{n,j}\delta_{\lambda_{n,j}}$ have mass $\kappa$
and support in $[0,\kappa]$. Hence their moments satisfy all three
quadratic-form inequalities in Section~\ref{sec:moments}, with
$M=\kappa$. Under \eqref{eq:hyp-power-target}, taking limits in
each finite quadratic form transfers those inequalities to the actual
$m_r$. The reconstruction theorem then supplies the calibrated seed.
\end{proof}

The fixed Hilbert space in this limit can be displayed directly.
The measures $\mu_n$ in the proof have uniformly bounded mass on
the compact interval $[0,\kappa]$, so they have weakly convergent
subsequences. Condition~\eqref{eq:hyp-power-target} fixes every
moment of any such limit. Polynomials are dense in the continuous
functions on this interval, so the limit measure is unique; call it
$\mu$. Then
\begin{equation}\label{eq:hyp-fixed-limit-seed}
\HH=L^2(\mu),\qquad (Kf)(t)=tf(t),\qquad v(t)=1
\end{equation}
satisfy $0\leq K\leq\kappa I$ and $\norm v^2=\kappa$.
The Taylor series of
$\ip{v}{(I+zK)^{-1}v}=\int(1+zt)^{-1}\dd\mu(t)$
has coefficients $(-1)^rm_r$, so it equals the actual $L$ near
zero. Thus the conditional limit constructs one fixed seed of the
type used in the main theorem.

For the squared pencils, the first two nonautomatic trace targets are
\begin{equation}\label{eq:hyp-first-trace-targets}
\lim_{n\to\infty}
 \frac{\Tr B_n^4}{(\Tr B_n^2)^2}=\frac{m_1}{\kappa^2},
\qquad
\lim_{n\to\infty}
 \frac{\Tr B_n^6}{(\Tr B_n^2)^3}=\frac{m_2}{\kappa^3}.
\end{equation}
These can serve as preliminary consistency tests for a proposed choice
of pencil parameters. They become an RH argument only when the
all-order statement \eqref{eq:hyp-power-target} is proved.

\begin{example}[Trace normalization does not give calibration]
Take $K_n=(\kappa/n)I_n$, which has nonnegative spectrum and trace
$\kappa$. Then
\[
Q_n(z)=(1+\kappa z/n)^n\longrightarrow e^{\kappa z},\qquad
\Tr(K_n^{r+1})=\kappa^{r+1}n^{-r}.
\]
This limit has no zeros, whereas the actual $\Theta$ has infinitely
many zeros. Thus positivity, growing dimension, exact normalization,
and locally uniform convergence do not by themselves identify the
limit with $\Theta$. The example is a general normalization test,
not an assertion that the project's particular pencils have this limit.
\end{example}

\begin{scopebox}
\small\textbf{The remaining Hyperbolic Zeta obligation.}
Choose an explicit family of admissible pencils and prove
\eqref{eq:hyp-coefficient-target}, equivalently
\eqref{eq:hyp-power-target}, for every degree. The constructions and
identities above establish a positive candidate family and the exact
consequence of such a limit theorem. They leave that limit theorem open.
The normalization uses the actual real constant $\kappa$, and the
pencil data use actual zeta values. Finite matrix dimension therefore
does not make this construction a finite raw checker witness.
The implication proved here is convergence of this family implies RH;
a converse for this particular family would require a separate
completeness theorem.
\end{scopebox}



\section[Finite matrices and witnesses]{Why a finite matrix is not the finite witness}\label{sec:finite}

\subsection{An obstruction to exact finite-dimensional calibration}

\begin{proposition}[Finite visible spectrum is impossible]\label{prop:finite-obstruction}
There is no exactly calibrated seed for the actual $\Theta$ whose
scalar spectral measure has finite support. In particular, no
finite-dimensional or finite-rank operator can supply an exact seed.
\end{proposition}
\begin{proof}
A finitely supported spectral measure makes
$F(z)=\sum_{j=1}^q c_j/(1+\lambda_jz)$ rational, with a possible
constant term when $\lambda_j=0$. Write $F=p/q_0$ for polynomials,
where $q_0$ is not identically zero. Exact interval calibration implies
that the entire function
\[
q_0(z)\Theta'(z)-p(z)\Theta(z)
\]
vanishes on an interval and therefore vanishes on all of $\C$.
At any zero $z_0$ of $\Theta$ with $q_0(z_0)\ne0$, the local
ODE $\Theta'=(p/q_0)\Theta$ contradicts the finite-order zero
argument of Theorem~\ref{thm:forward}. Thus every zero of $\Theta$
would have to be among the finitely many roots of $q_0$.
But $\xi$ has infinitely many distinct zeros, and the two-to-one
square map can identify only opposite pairs. Hence $\Theta$ has
infinitely many distinct zeros, a contradiction.

If $K$ is finite rank, its nonzero spectral values form a finite set;
its kernel can contribute only the additional spectral point zero.
Its scalar spectral measure is therefore finitely supported, so the
same argument applies.
\end{proof}

This obstruction does not prohibit a finite description of an
infinite-dimensional operator or a finite certificate for an infinite
analytic argument. It rules out the narrower proposal that a finite
positive matrix by itself can reproduce the full actual logarithmic
derivative exactly on an interval.

\subsection{The two missing checker theorems}

The project's requested finite endpoint is
\begin{equation}\label{eq:finite-goal}
\bigl(\exists W:\mathsf{FiniteWitness},\ \mathsf{check}(W)=\mathsf{true}\bigr)
\quad\Longleftrightarrow\quad\RH.
\end{equation}
The witness must consist of finite raw data and the checker must
terminate. The analytic seed defined in this paper is not such a
witness type: it quantifies over a Hilbert space, an operator, a vector,
and an exact interval identity.

A route through Theorem~\ref{thm:main} would require the following
separate statements for one specified checker:
\begin{align}
\mathsf{check}(W)=\mathsf{true}&\ \Longrightarrow\ \HSeed,
&&\text{for every finite }W;\label{eq:soundness}\\
\RH&\ \Longrightarrow\ \exists W,\ \mathsf{check}(W)=\mathsf{true}.
&&\label{eq:completeness}
\end{align}
Equation~\eqref{eq:soundness} is actual-function soundness. It must
derive calibration, not receive an unverified calibration proof as a
field of the purported raw witness. Equation~\eqref{eq:completeness}
is finite-language completeness. Constructing a countable spectrum
from all zeros under RH proves an analytic existence statement; it
does not extract an accepted finite code in an independently fixed
language.

Finite moment matrices, finite Gram matrices and finite truncations
are legitimate components of a possible certificate design. To yield
\eqref{eq:soundness}, they need a theorem extending the verified finite
information to the uniform analytic assertion. The all-order moment
criterion in Corollary~\ref{cor:moment-criterion} specifies one precise
target for such an extension. No such finite reduction is established
by the present note.

\subsection{Finite jets cannot replace a functional identity}

There is a direct mathematical reason for the distinction. Fix a
real anchor $a_*>0$ and an integer $N\geq0$. The functions
\[
L(z)\quad\text{and}\quad L(z)+\varepsilon(z-a_*)^{N+1}
\]
have identical derivatives through order $N$ at $a_*$ for every
$\varepsilon$, but are different analytic functions when
$\varepsilon\ne0$. Finitely many samples can likewise be preserved by
adding a nonzero polynomial vanishing at all sample points.
The identity theorem requires a set with an accumulation point in
the analytic domain, or equivalent exact all-order germ information.
A finite initial jet is neither of these.

\section[Repository and evidence]{Connection with the repository and proof obligations}\label{sec:repository}

This note was prepared against the public repository
\texttt{randogelist/Riemann-Hypothesis-Final} at the source revision
\begin{center}\small\reposha.\end{center}
The inspected source records an analytic slit-plane admission
\cite{repo-ode}:
\begin{equation}\label{eq:repo-admission}
\exists A\text{ holomorphic on }\C\setminus(-\infty,0],\qquad
\Theta'+A\Theta=0\text{ there}.
\end{equation}
The source declaration
\nolinkurl{SoloboloAdmission.slitODEAdmission_iff_RH}
has literal target \texttt{\_root\_.RiemannHypothesis}.
The resolvent transfer in this paper supplies the coefficient
$A=-F$ on that domain, because the smaller slit plane is contained
in $D_M$. The sign is negative: the paper writes $\Theta'=F\Theta$,
whereas the source admission writes $\Theta'+A\Theta=0$.

The source zero-geometry declaration
\nolinkurl{AngularMaster.native_RH_iff_theta_zeros_nonpositive}
matches Proposition~\ref{prop:geometry}, allowing nonpositive real
zeros before origin nonvanishing is supplied \cite{repo-geometry}.
The repository's positive-product assessment uses the normalization
$\sum_n\lambda_n/(1+a\lambda_n)$ with
$\lambda_n=\gamma_n^{-2}$, consistent with
\eqref{eq:diagonal-calibration} \cite{repo-product}.

The particular legacy archive name
\nolinkurl{RH_SOLOBOLO_SEED_v0_1_0.zip} and the named
\texttt{NativeOperatorSeed} sources were not recovered in the inspected
repository tree or targeted searches. Accordingly, Definitions
\ref{def:hseed} and \ref{def:dseed} are explicit mathematical
formulations supplied by this paper. No claim is made that their
fields are textually identical to an unrecovered legacy declaration.

\clearpage
\begin{longtable}{@{}>{\raggedright\arraybackslash}p{0.19\textwidth}>{\raggedright\arraybackslash}p{0.47\textwidth}>{\raggedright\arraybackslash}p{0.27\textwidth}@{}}
\toprule
\textbf{Layer}&\textbf{Statement or construction}&\textbf{Evidence in this note}\\
\midrule
\endhead
Resolvent algebra&Norm-square identity and all-parameter Gram formula&On-paper proofs, Sections~3--4\\
\addlinespace
Analytic transfer&Exact normalized or differential seed implies actual RH&On-paper proofs, Sections~5--6\\
\addlinespace
Converse&RH gives the positive countable diagonal seed&On-paper proof from classical factorization, Section~7\\
\addlinespace
Moment route&All-order positivity with one support bound gives a seed&On-paper reconstruction, Section~8\\
\addlinespace
Hyperbolic zeta&Exact model integrals, actual-jet identities, and a conditional determinant-limit transfer&On-paper proofs, Section~\ref{sec:hyperbolic}; convergence to actual coefficients remains open\\
\addlinespace
Legacy seed&Exact \texttt{SeedOn}/\texttt{HasSeed} implementation&Not recovered or certified here\\
\addlinespace
Lean verification&Compilation and exact endpoint axiom audit&Not run; no new formal certification\\
\addlinespace
Finite checker&Actual soundness, finite completeness, accepted witness&Open; not constructed here\\
\bottomrule
\end{longtable}

The repository's inspected status records native build and axiom audit
as \texttt{NOT\_RUN}, actual calibration as \texttt{OPEN}, and
finite-witness soundness and completeness as \texttt{OPEN}
\cite{repo-status}. Publication of this paper does not change those
statuses. A later formalization must connect the exact definitions,
complex derivative, domains, spectral constructions and zero geometry
to the actual native declarations; familiar theorem names are not a
substitute for that transport.

\section{Conclusion}\label{sec:conclusion}

The recovered Hilbert-vector anchor has a complete and coherent
on-paper analytic formulation. Its central scalar identity is
\[
\frac{\mathrm d}{\mathrm da}\left(a\frac{\Theta'(a)}{\Theta(a)}\right)
=\norm{(I+aK)^{-1}v}^2,
\]
where $K$ and $v$ are fixed, $K$ is bounded and nonnegative
self-adjoint, and the identity is exact on a positive interval.
The proof proceeds through a holomorphic resolvent, an explicitly
removed integration constant, and a division-free analytic ODE.
The zero-order contradiction then yields the critical-line statement.

Conversely RH supplies a diagonal seed through the genus-zero
product of the entire square-descended function. This gives the
analytic equivalence \eqref{eq:main}, the Gram identity, and the
normalization \eqref{eq:trace-norm}. It does not construct the seed
without RH. The unresolved constructive task is to derive one of
these exact analytic inputs from independently verifiable data;
the further finite-checker equivalence requires both
\eqref{eq:soundness} and \eqref{eq:completeness}.


Section~\ref{sec:hyperbolic} gives a precise role for the Hyperbolic
Zeta identities. The Bose weights construct positive bounded model
resolvents and exact odd-zeta formulas; their asymptotics exclude a
direct identification with the actual $L$. A separate positive
hyperbolic-tangent integral determines actual $\xi$ jets, and the
Newton recurrence converts those jets to the required logarithmic
moments. Squaring the project's signed pencils supplies positive
determinant models. Theorem~\ref{thm:hyp-coefficient-transfer} proves
that calibration of all their coefficients, or all their trace moments,
would construct the fixed seed and imply RH. That all-order limit
theorem, and a completeness theorem for this particular family,
remain open.

\appendix
\section[Real Hilbert spaces]{Real Hilbert spaces and the holomorphic convention}\label{app:real}

The recovered anchor may be expressed on a real Hilbert space.
That does not obstruct the complex-analytic proof. If $H$ is real,
form $H_\C=H\oplus iH$ with
\begin{align*}
\ip{x+iy}{u+iw}_{\C}
&=(x,u)_H+(y,w)_H\\
&\quad+i\bigl((x,w)_H-(y,u)_H\bigr).
\end{align*}
This inner product is conjugate-linear in the first slot and linear
in the second. A real bounded self-adjoint nonnegative $K$ extends
by $K_\C(x+iy)=Kx+iKy$. Its norm is unchanged, because
\[
\norm{K_\C(x+iy)}^2=\norm{Kx}^2+\norm{Ky}^2
\leq\norm K^2(\norm x^2+\norm y^2),
\]
and restriction to real vectors gives the reverse norm inequality.
Self-adjointness follows by expanding the inner product. Positivity
follows from
\[
\ip{x+iy}{K_\C(x+iy)}=(x,Kx)_H+(y,Ky)_H\geq0,
\]
since the imaginary cross terms cancel by symmetry.
For real $a\geq0$, the inverse extends the real inverse, so every
real scalar calibration and norm identity is preserved. Thus the
complex theorem applies to the real Hilbert version as well.

If an alternative convention makes the inner product linear in the
first slot, the scalar holomorphic expression should instead be
$\ip{R_zv}{v}$. With the convention fixed in this paper, it is
$\ip v{R_zv}$. Mixing these conventions would turn a holomorphic
function into an anti-holomorphic one and invalidate the identity-theorem
step.

\section[Proof dependencies]{A compact proof dependency ledger}\label{app:ledger}

The forward implication has the following dependencies, each proved
or cited at its point of use:
\begin{enumerate}[label=\textbf{F\arabic*.},leftmargin=2.8em]
\item The actual $\Theta$ is entire and $\Theta(0)>0$.
\item Fixed bounded positive $K$ gives a holomorphic resolvent on
      $D_M$, a connected domain containing zero and the positive ray.
\item Exact interval calibration gives $\Theta'=F\Theta$ by the
      identity theorem. For differential calibration, first remove
      $C$ using $Q(0)=-C\Theta(0)$.
\item A nonzero entire solution of that ODE has no zero where its
      coefficient is holomorphic, by the finite-order contradiction.
\item Negative-real-axis confinement of $\Theta$ zeros is equivalent
      to the critical-line condition for actual $\xi$ zeros.
\end{enumerate}
Only the converse uses RH:
\begin{enumerate}[label=\textbf{R\arabic*.},leftmargin=2.8em]
\item RH places the zeros of $\Theta$ at $-\gamma_n^2$.
\item The order bound and Hadamard factorization give
      $\Theta/\Theta(0)=\prod_n(1+z/\gamma_n^2)$ and summability.
\item Uniformly convergent logarithmic differentiation gives
      $L(a)=\sum_n(a+\gamma_n^2)^{-1}$.
\item The explicit diagonal $K$ and square-summable $v$ realize this
      scalar function and hence both seed conditions.
\end{enumerate}
The all-order moment criterion reconstructs the forward input through
Theorem~\ref{thm:gns}. It does not remove an infinite quantifier by
itself. Neither direction in this ledger uses a finite accepted
witness, and neither is reported as a Lean-verified endpoint.

\begin{thebibliography}{99}
\small
\bibitem{dlmf-reflection}
NIST Digital Library of Mathematical Functions,
\emph{Section 25.4: Reflection Formulas}, especially Eqs.~25.4.3--25.4.4.
\url{https://dlmf.nist.gov/25.4}. Accessed 27 October 2026.

\bibitem{dlmf-series}
NIST Digital Library of Mathematical Functions,
\emph{Section 25.2: Definition and Expansions}, especially Eq.~25.2.3.
\url{https://dlmf.nist.gov/25.2}. Accessed 27 October 2026.

\bibitem{dlmf-zeros}
NIST Digital Library of Mathematical Functions,
\emph{Section 25.10: Zeros}, subsection (i).
\url{https://dlmf.nist.gov/25.10}. Accessed 27 October 2026.

\bibitem{garrett}
Paul Garrett, \emph{Riemann-Hadamard product for $\zeta(s)$},
24 October 2021, pp.~1--2 (growth of the completed entire function).
\url{https://www-users.cse.umn.edu/~garrett/m/number_theory/riemann_hadamard_product.pdf}.

\bibitem{mcmullen}
C. McMullen, \emph{Advanced Complex Analysis}, Harvard Math 213a
course notes, 22 January 2017. Theorems~3.7 and 3.11,
printed pp.~43--44 (Hadamard factorization and zero summability).
\url{https://people.math.harvard.edu/~ctm/home/text/class/harvard/213a/10/html/home/course/course.pdf}.

\bibitem{williams}
Dana P. Williams, \emph{Lecture Notes on the Spectral Theorem},
10 August 2018, Sections~2, 4 and 5 (bounded spectral calculus and
spectral measures).
\url{https://math.dartmouth.edu/~dana/bookspapers/ln-spec-thm.pdf}.

\bibitem{repo-ode}
Riemann-Hypothesis-Final project, \texttt{SoloboloAdmission/NativeODE.lean},
source revision \texttt{362d9f8}. Exact source of
\texttt{slitODEAdmission\_iff\_RH}; native audit not run in this note.
\href{https://github.com/randogelist/Riemann-Hypothesis-Final/blob/362d9f8aeaaf9479cc27ceb6b63bdb7456267990/SoloboloAdmission/NativeODE.lean}{Pinned repository source}.

\bibitem{repo-geometry}
Riemann-Hypothesis-Final project, \texttt{AngularMaster/ZeroGeometry.lean},
source revision \texttt{362d9f8}.
\href{https://github.com/randogelist/Riemann-Hypothesis-Final/blob/362d9f8aeaaf9479cc27ceb6b63bdb7456267990/AngularMaster/ZeroGeometry.lean}{Pinned repository source}.

\bibitem{repo-product}
Riemann-Hypothesis-Final project,
\emph{RH positive-product assessment}, 25 October 2026,
source revision \texttt{362d9f8}.
\href{https://github.com/randogelist/Riemann-Hypothesis-Final/blob/362d9f8aeaaf9479cc27ceb6b63bdb7456267990/docs/research/2026-10-25/rh-positive-product-assessment.md}{Pinned research assessment}.

\bibitem{repo-status}
Riemann-Hypothesis-Final project, \texttt{provenance/status.json},
\texttt{docs/BUILD.md}, and \texttt{docs/FRONTIER.md},
source revision \texttt{362d9f8}. Source evidence and open finite
calibration obligations.
\href{https://github.com/randogelist/Riemann-Hypothesis-Final/blob/362d9f8aeaaf9479cc27ceb6b63bdb7456267990/provenance/status.json}{Pinned status record}.
\bibitem{hyp-manuscript}
Prince Anand, \emph{Spectral Theory and Diophantine Bounds for Hankel
Determinants of Odd Zeta Values}, supplied project manuscript,
printed pp.~7--9, especially Eq.~(3.1) and Theorem~3.1.
Tracked in \texttt{randogelist/zetahyperbolic-axiom}, revision
\texttt{b56c8102d25b931a601992c93bf30b19cc7ee103}.
\href{https://github.com/randogelist/zetahyperbolic-axiom/tree/b56c8102d25b931a601992c93bf30b19cc7ee103}{Pinned manuscript repository}.
Only the identities identified in Section~\ref{sec:hyperbolic} are
used here; other manuscript endpoints are outside this note's scope.

\bibitem{hyp-weight}
zetahyperbolic-axiom project,
\texttt{Theorem\_3\_1\_Hyperbolic\_Weight.lean},
revision \texttt{b56c810}, declarations \texttt{density} and
\texttt{density\_pos}. Source inspection; no fresh native audit.
\href{https://github.com/randogelist/zetahyperbolic-axiom/blob/b56c8102d25b931a601992c93bf30b19cc7ee103/src/Theorem_3_1_Hyperbolic_Weight.lean}{Pinned weight source}.

\bibitem{hyp-moments}
zetahyperbolic-axiom project,
\texttt{Theorem\_3\_1\_Hyperbolic\_Moments.lean},
revision \texttt{b56c810}, declaration
\texttt{integral\_monomial\_density}.
\href{https://github.com/randogelist/zetahyperbolic-axiom/blob/b56c8102d25b931a601992c93bf30b19cc7ee103/src/Theorem_3_1_Hyperbolic_Moments.lean}{Pinned moment source}.

\bibitem{hyp-pole}
zetahyperbolic-axiom project,
\texttt{Theorem\_3\_1\_Pole\_Integral\_Identity.lean},
revision \texttt{b56c810}, declaration \texttt{poleIntegralIdentity}.
Its integer-parameter statement is extended on paper to real positive
parameters in Proposition~\ref{hyp-hurwitz-prop}.
\href{https://github.com/randogelist/zetahyperbolic-axiom/blob/b56c8102d25b931a601992c93bf30b19cc7ee103/src/Theorem_3_1_Pole_Integral_Identity.lean}{Pinned pole-integral source}.

\bibitem{hyp-normalization}
zetahyperbolic-axiom project,
\texttt{LeanPaper/ContinuousDensityNormalization.lean},
revision \texttt{b56c810}, declarations
\texttt{correctedDensity\_odd} and \texttt{printedDensity\_odd}.
\href{https://github.com/randogelist/zetahyperbolic-axiom/blob/b56c8102d25b931a601992c93bf30b19cc7ee103/src/LeanPaper/ContinuousDensityNormalization.lean}{Pinned normalization correction}.

\bibitem{hyp-gram}
zetahyperbolic-axiom project,
\texttt{LeanPaper/ContinuousGram.lean}, revision \texttt{b56c810},
declarations \texttt{baseMatrix\_entry\_exact},
\texttt{baseMatrix\_posDef}, and \texttt{pencil\_odd}.
\href{https://github.com/randogelist/zetahyperbolic-axiom/blob/b56c8102d25b931a601992c93bf30b19cc7ee103/src/LeanPaper/ContinuousGram.lean}{Pinned Gram source}.

\bibitem{hyp-signs}
zetahyperbolic-axiom project,
\texttt{LeanPaper/PoleWeightSigns.lean}, revision \texttt{b56c810},
declarations \texttt{poleWeight\_sign} and
\texttt{gamma\_has\_both\_signs}.
\href{https://github.com/randogelist/zetahyperbolic-axiom/blob/b56c8102d25b931a601992c93bf30b19cc7ee103/src/LeanPaper/PoleWeightSigns.lean}{Pinned sign source}.
See also \texttt{snditeration}, revision
\texttt{35218ba6ee152774ff4e1fb778e530c1752acfd6},
\texttt{Zetahyperbolic/Section3/SpectralInertia.lean},
\texttt{actual\_inertia\_formula}.

\bibitem{dlmf-binet}
NIST Digital Library of Mathematical Functions,
\emph{Section 5.9: Integral Representations}, Eq.~5.9.15
(Binet's digamma formula).
\url{https://dlmf.nist.gov/5.9}. Accessed 27 October 2026.

\bibitem{dlmf-hurwitz}
NIST Digital Library of Mathematical Functions,
\emph{Section 25.11: Hurwitz Zeta Function},
Eqs.~25.11.1, 25.11.3, 25.11.12, and 25.11.17.
\url{https://dlmf.nist.gov/25.11}. Accessed 27 October 2026.

\bibitem{dlmf-mangoldt}
NIST Digital Library of Mathematical Functions,
\emph{Section 27.4: Euler Products and Dirichlet Series}, Eq.~27.4.12.
\url{https://dlmf.nist.gov/27.4}. Accessed 27 October 2026.

\bibitem{dlmf-fermi}
NIST Digital Library of Mathematical Functions,
\emph{Section 25.5: Integral Representations}, Eq.~25.5.3.
\url{https://dlmf.nist.gov/25.5}. Accessed 27 October 2026.

\bibitem{dlmf-duplication}
NIST Digital Library of Mathematical Functions,
\emph{Section 5.5: Functional Relations}, Eq.~5.5.5
(gamma duplication).
\url{https://dlmf.nist.gov/5.5}. Accessed 27 October 2026.

\bibitem{csordas-kernel}
George Csordas, \emph{Fourier transforms of positive definite kernels
and the Riemann $\xi$-function}, arXiv:1309.0055v2,
21 February 2014, Section~4, Eqs.~(4.1)--(4.2), printed pp.~9--10.
The Fourier-variable $\xi$ in that source is distinguished from
the completion $\xi(s)$ used here by the rescaling in
\eqref{hyp-eq:riemann-cosh}.
\url{https://arxiv.org/abs/1309.0055}.

\end{thebibliography}
\end{document}
